\chapter{Overhead}
\label{chap:overhead}

\begin{chapteressay}
Overhead is not automatically corruption. Cheap histories of science can
make every institution look like a conspiracy and every instrument with a
budget look guilty. That is not the argument here. A large detector, a
registry, a standard laboratory, a calibration chain, a review board, a
journal, a statistical center, a field station, and a training program all
cost something. Without that cost, many modern records would not exist. The
danger is subtler: cost can learn to speak as if it were evidence.

Staged values have become portable. Portability is never free. If a value
travels, someone maintains the standard, teaches the procedure, preserves
the archive, pays the technicians, reviews the claim, compares the
instruments, stores the data, publishes the result, and decides which
finite record stands for the event stream the world actually produced.
These are not embarrassments around science. They are part of the
historical machinery by which science becomes public.

The machinery is dangerous because it has its own incentives. A journal
wants news. A laboratory wants continuation. A field wants coherence. A
funding program wants deliverables. A government wants policy. A company
wants an approved product. A public wants a story. None of those desires is
evidence, but each can attach itself to evidence and begin to carry it
farther than the record allows. Overhead is the zone where support and
amplification become hard to distinguish.

Overhead carries \lean{LOAD}, \lean{FINITE_ELEPHANT}, \lean{BULLSHIT},
\lean{PROPAGANDA}, and \lean{ACOLYTE}. The satirical names matter because
polite names let the failure hide. \lean{LOAD} is the real burden of
keeping a record alive. \lean{FINITE_ELEPHANT} is the honest confession
that the record cannot carry everything. \lean{BULLSHIT} is not mere
falsehood; it is a claim whose social performance is indifferent to whether
the route holds. \lean{PROPAGANDA} is overhead that has learned to defend a
sentence by means other than evidence. \lean{ACOLYTE} is the person or
institution that repeats the sentence because the carrier rewards
repetition.

Overhead keeps two refusals alive at once. First, it refuses cynicism: the
existence of overhead cannot make the record fake. The LHC trigger is
overhead and is honest precisely because it publishes its finite selection.
Clinical-trial registries are overhead and are honest because they keep
unflattering absences from disappearing. Standards bodies are overhead and
are honest when they keep local instruments answerable to public
operations.

Second, it refuses credulity: overhead may be necessary, but necessity
cannot sanctify it. N-rays had apparatus, prestige, publication, and
repeated performances. Cold fusion had laboratories, announcements,
journals, and institutional hopes. Lysenkoism had state authority,
educational machinery, and acolytes. Piltdown had museums and experts. None
of that made the record hold. The carrier can be large and still carry a
false claim.

The distinction is not moral purity. Honest overhead can still be
political, expensive, selective, and fallible. Bad overhead can still
contain sincere people and real instruments. The operational question is
whether the overhead remains answerable to the staged route. Outsiders can
see what was selected, what was lost, what was corrected, what failed to
replicate, what uncertainty remains, and where the claim stops --- or the
overhead converts those questions into disloyalty, ignorance, heresy, or
public-relations risk.

That is why overhead follows staging. A claim with no staged value lacks
the public weight to attract this kind of support. Once a value travels,
however, it becomes worth defending, funding, teaching, and weaponizing.
The very success of experimental science creates the social conditions in
which experimental authority can be imitated.

The cases move between honest and dishonest overhead. N-rays and cold
fusion show expectation and announcement outrunning the record. Lavoisier
and Lysenko show that institutions can either preserve an experimental
refusal or punish it. The LHC trigger shows finite loss made public as
method. Climate consensus appears as a boundary problem rather than a
slogan: measured warming, attribution, projections, policy choice, and
public ritual stay distinguished rather than collapsing into one
undifferentiated authority. Big observatories, citation cascades,
replication projects, trial registries, polywater, Piltdown, and Shewhart
charts keep the argument grounded in the ordinary ways claims are carried
or corrected.

Overhead is uncomfortable because it talks about science where science is
socially embodied. That is also why it is necessary. The universe is not
publicly inferable through pure marks alone. It is inferable through
records carried by people, instruments, committees, institutions,
archives, and habits. Those carriers are part of the truth's route. They
are not allowed to become the truth.

Overhead also has a time signature. It often arrives after a value has
begun to travel, but it can also arrive before the value is ready. A
laboratory press office can make a premature claim public. A journal
culture can teach young researchers which shapes of result are
publishable. A funding system can make some questions easy to ask and
others nearly invisible. A standards office can keep a record alive for a
century. Overhead is the historical weather in which experimental claims
either keep their routes or lose them.

The practical test is whether the carrier remains reversible. Tracing the
public sentence back to the apparatus, protocol, selection, uncertainty,
and refusal is possible, or the carrier has become a substitute world.
That substitute world may be charming, prestigious, terrifying,
profitable, or morally satisfying. None of those qualities make it
evidence.

Celebrating experiments that succeeded while mocking experiments that
failed misses the point. The point is what kinds of carrying structures
let success and failure become visible. A failed replication can be a triumph of overhead
when the system lets the refusal be published. A famous claim can be a
failure of overhead when the system protects it from the conditions that
would make it answerable.

Three overhead failures recur. The first is substitution: the carrier
performs the work that belongs to the record. Prestige substitutes for signal, announcement for replication,
consensus for boundary, citation for root evidence. The second is
insulation: the carrier protects a claim from the refusals that would
discipline it. Failed replications become incompetence, critics become
enemies, uncertainty becomes disloyalty, missing results become invisible.
The third is capture: the claim becomes useful to a social structure that
now has reasons to keep the claim alive apart from the evidence.

The honest forms have their own three virtues. First, declaration: the
carrier says what it adds. A trigger declares selection; a registry
declares a trial; a control chart declares limits; a standards lab
declares uncertainty. Second, reversibility: the public sentence can be
traced back toward the route. A citation points to the root, an archive
carries provenance, an assessment separates measured record from
projection. Third, corrigibility: the carrier can be corrected by the
record it carries. These virtues are not decorations. They are how
overhead remains load-bearing.

Two easy moods are worth resisting. The first is suspicion, in which every
institution is assumed to be lying because institutions have interests.
The second is deference, in which every institution is assumed to be
right because institutions have expertise. Both moods are lazy. Experimental history is harder and better:
institutions are necessary carriers that remain under discipline only
when the route is kept inspectable.

The satirical names sit in serious argument because they make the failure
types memorable without becoming vague insults. \lean{BULLSHIT} is not
``a claim I dislike.'' It is discourse detached from the record's route.
\lean{PROPAGANDA} is not ``politics near science.'' It is a carrier
defending a sentence by means that bypass evidence. \lean{ACOLYTE} is not
``a person with consensus views.'' It is repetition rewarded by the carrier
after the route has been subordinated. The names are sharp only when the
diagnosis is sharp.

The diagnostic is practical. Ask first what the carrier adds: memory,
calibration, distribution, review, storage, comparison, or selection.
Those additions may be necessary. Ask second what the carrier hides:
missing trials, discarded events, weak roots, changed endpoints,
uncertainty, or social pressure. Ask third what happens when the record
says no: the carrier publishes the refusal, reinterprets it honestly,
buries it, punishes it, or turns it into a loyalty test. The third
question is usually the decisive one.

This diagnostic also explains why overhead cannot be audited by tone. A
calm institutional voice can carry propaganda. A furious critic can carry
a real refusal. A boring registry can protect truth. A beautiful image
can distract from calibration. A prestigious consensus can summarize
evidence, and it can also overrun boundaries. A fringe objection can be
bad faith, and it can also point to a real gap. Tone is not route.

Overhead also cannot be audited by novelty. New claims can be true or
false. Established claims can be well grounded or merely inherited. The
question is always whether the carrier remains attached to the operations
that made the claim answerable. N-rays and clinical registries belong to the same problem: the public route
of a claim after people begin to carry it.

The ambition is therefore modest and severe. It cannot purify science of
social life; it can make social life answerable. A budget, institution,
consensus, citation, or announcement has a job and a limit. The carrier
shows the route, or the carrier is refused in place of the route.

A realistic history of experiment cannot stop at the apparatus drawing.
It includes the meeting where the result was announced, the archive where
the record was stored, the citation that carried the claim into another
field, the committee that selected the event, the registry that kept a
trial from disappearing, and the classroom where a student learned which
doubts were acceptable. These are not anecdotes outside science. They are
the social surfaces along which experimental claims move.

Not every social surface matters equally. The balance matters more than
Lavoisier's reputation when the question is phlogiston. The trigger rule
matters more than the collaboration logo when the question is an event
excess. The trial endpoint matters more than the sponsor's slogan when
the question is efficacy. Overhead analysis is the discipline of ranking
such surfaces by their relation to the route.

That ranking keeps overhead analysis from becoming gossip. The point is
not scandal but the place where the carrier changes the evidential status
of the claim. When the carrier preserves route, the case credits it. When
the carrier hides or replaces route, the case names the failure. The same
institution can act either way in different moments, which is why the
diagnostic stays awake.

Staying awake means refusing easy moral sorting. N-rays involved sincere
observers. Cold fusion involved real apparatus. Climate assessment involves
real evidence. Big observatories involve real discovery. Citation cascades
and replication failures often involve ordinary scholars acting on trained
incentives. The gauge names failure without requiring villainy.

It also names virtue without requiring purity. Honest triggers, registries,
control charts, archives, and corrections come from fallible people building
forms that make fallibility less decisive. That is the hope: design against
self-deception.

This design is not glamorous. A culture that only celebrates breakthroughs
underfunds the forms that keep breakthroughs honest. The registry, archive, control
limit, correction notice, calibration log, and trigger table may look like
supporting paperwork. In experimental history, they are part of the
apparatus by which public truth survives its own popularity.
\end{chapteressay}

\section{N-Rays: Expectation Masquerading as Signal}
\episodecoord{1903-1906 | Overhead :: BULLSHIT -> PROPAGANDA}

N-rays were more than a failed residue. Their failure was not only optical
or experimental. It was social. The
claimed rays were seen by trained people in prestigious settings, written up
in publications, repeated in variations, and made to carry the authority of an
eminent laboratory. The record carried more than a mistaken observation. It
carried overhead that helped a mistaken observation behave like a field.

René Blondlot's reports described a new radiation detected through small
changes in luminosity. The apparatus could look delicate and plausible: a
source, a prism or screen, a detector surface, an observer judging a faint
change. The problem was that the crucial mark lived too close to expectation.
The observer was asked to say whether a barely visible change had occurred. When
the signal depends on trained anticipation at the edge of perception, social
force can enter the measurement without leaving a mechanical trace.

The point is not that the observers were stupid. The staging was
under-governed for the claim being made. The eye was being asked to carry too
much. The apparatus failed to separate source behavior from the observer's
readiness to see source behavior. The experimental route left a
gap, and the gap was filled by expectation, reputation, and the desire to
participate in a discovery.

Robert W. Wood's intervention is famous because he disturbed the apparatus and
the claimed observations continued. In the cleanest telling, he secretly
removed a prism and the observers still reported the effect. The story cannot
be polished beyond what the record supports, but the structural lesson is
clear: a staged refusal was introduced into the overhead. If the apparatus no
longer contained the supposed element required for the effect, continued
observation of the effect exposed the observation as expectation-laden.

This is the difference between residue and overhead. As residue, N-rays fail
the replication gate. As overhead, they show how a failing record can be
carried. Laboratories repeated procedures, journals printed reports,
colleagues looked for effects, and national pride supplied warmth. None of
those forces created the rays. They created a climate in which the claim could
feel experimentally alive after the route had failed to earn the signal.

\lean{BULLSHIT} is the precise word here because the claim's public force
became increasingly detached from whether the phenomenon was there. The social
performance of seeing mattered more than the disciplined route by which seeing
would be made public. The effect's defenders were not simply lying in the
ordinary sense. They were participating in a discourse whose relation to the
record had loosened. The claim could be repeated, elaborated, and defended
without acquiring the controls that would make it answerable.

The transition toward \lean{PROPAGANDA} begins when the claim recruits a
community identity. A discovery can become "ours." A criticism can become an
attack. A refusal can become an insult rather than a correction. The
experimental question then changes character. Instead of asking whether the
route holds, the community asks who is willing to see, who is sympathetic,
who understands the delicate method, who belongs. That is the acolyte path
before the name \lean{ACOLYTE} fully appears.

N-rays are therefore not merely a joke about French physics or visual
suggestion. They are a warning about low-signal experimental cultures. When
the mark is faint, the status of the observer becomes tempting. The answer
is not to ban faint marks; much real science begins in weak signals. The
answer is to make the route severe enough that expectation cannot carry the
claim by itself. Blind controls, mechanical records, independent replication, and adversarial
disturbances exist because observers are human.

The case also guards against a common misuse of failure stories. The
lesson is not that surprising effects are probably false. It is that the
more surprising and delicate the effect, the more the route protects itself
against the social energy that surprise attracts. New phenomena are allowed
to be difficult. They are not allowed to substitute prestige for controls.

N-rays also show why a failed field can be historically useful. A scientific
community learns from the embarrassment by sharpening the difference between
private acuity and public detection. The observer's eye remains valuable; the
eye is put under discipline. Apparatus is disturbed, blinded, rearranged, and
made to answer in ways independent of a person's desire to report the
expected result.

The case is especially important because it keeps overhead from
being restricted to governments or giant machines. Overhead can gather around
a tabletop, a dim screen, a respected name, and a room full of expectation.
It can be small, elegant, and sincere. The gauge therefore detects structure
rather than size. The failure signature is not expense. The failure
signature is a carrier that keeps moving after the experimental route has
fallen away.

Delicacy is not vagueness. Delicate experiments are allowed to require
trained hands and eyes. The problem with N-rays was not that the report
asked observers to attend carefully. Many real effects demand careful
attention. The problem was that the care could not close the route. A
trained report remained too exposed to suggestion because the apparatus
could not force the signal to answer independently enough. The experiment made
expertise heavier than the mark could bear.

The social reward for seeing then becomes part of the apparatus. A young
researcher in a prestigious laboratory may know, without anyone saying it
crudely, which answer is welcome. A visitor may want to be polite. A skeptic
may be told they lack the necessary sensitivity. A defender may interpret
failure as clumsiness rather than refusal. These pressures are usually soft,
and soft pressures are exactly what low-signal experiments are designed to
survive.

Blinded or adversarial intervention is not an insult. It is a gift to the
record. A scientist who believes a faint effect is real welcomes a procedure that
protects the effect from the scientist's own desire. Wood's
intervention was severe because it separated the social act of reporting from
the physical condition supposedly required for the report. Once those came
apart, the overhead was exposed.

The case also shows how quickly a field can generate elaboration before it
generates discipline. New sources, new materials, new biological or optical
effects, and new explanatory sketches can proliferate around a weak root. Each
addition makes the claim feel richer. But richness is not custody. A weak
root with many branches remains weak until the root is made public under
control.

N-rays therefore teach a lesson that returns in modern forms whenever a
signal sits close to the limits of perception, preprocessing, or statistical
selection. The exact apparatus changes, but the danger remains: the carrier
can learn the expected answer. The only defense is a route that can disappoint
the carrier.

The N-ray story is therefore not the tale of one clever debunker saving
science from foolish believers. That version is too comforting. The
better lesson is systemic. Any community that rewards a difficult perception
can train perception. Any laboratory that makes membership depend on seeing
what the master sees can turn the apparatus into a social test. Any field
that treats criticism as crudity can confuse delicacy with immunity. Wood's
role matters because he forced a private delicacy into public interruption,
but the larger question is why interruption had become necessary.

The N-ray episode also exposes the difference between expertise and authority.
Expertise can help an observer notice a real weak signal. Authority can make a
weak signal harder to challenge. The same trained person may carry both. A
healthy experimental culture builds procedures that let expertise work while
limiting authority's ability to decide the result. Controls are not
anti-expert. They protect expertise from being confused with status.

The failure had a pedagogical afterlife. Later accounts associated N-rays
with suggestion, national rivalry, and the need for blind or mechanical
checks. That afterlife is useful only if it is not used as a sneer. The
question is not whether "they" were gullible and "we" are modern. The question
is where today's instruments still ask human or algorithmic expectation to
carry too much. A faint image, a borderline p-value, a manually selected
feature, a processed detection, or a classifier's output can all become an
N-ray if the route cannot surprise the expectation.

The word \lean{PROPAGANDA} enters gently here. N-rays required no ministry.
They show a proto-propaganda form: a claim becomes attached to group pride
and distributed as a sign of sophistication. Once that happens, refusal
fights not only an observation but a role. The person who refuses to confirm
is no longer simply reporting a negative result. They are refusing a place
in the community's self-description.

That is the soft beginning of many hard failures. The acolyte first learns to
see. Then to repeat. Then to defend. The record may still interrupt the chain,
but each social step makes interruption more costly. N-rays remain valuable
because the interruption came before the carrier hardened into something more
durable.

\begin{phenom}{The N-Ray Overhead Effect~\cite{blondlot1904,wood1904}}
\label{ph:n-ray-overhead}

\PhStatement
Expectation becomes overhead when the social performance of detecting a signal
is allowed to substitute for a public route to the signal.

\PhOrigin
Blondlot reported N-rays through delicate luminosity observations. Wood's
critical intervention and failed replication exposed the claimed effect as
unreliable.

\PhObservation
The record includes faint visual judgments, apparatus arrangements, published
reports, laboratory prestige, and critical disturbance of the setup.

\PhConstraint
Low-signal claims are staged so that expectation, reputation, and observer
desire cannot carry the mark in place of the apparatus.

\PhConsequence
The failure is not that observers can be fooled. The failure is that overhead
can make being fooled look like membership in a discovery.
\end{phenom}

The bridge from N-rays to cold fusion is the bridge from perception to
announcement. In N-rays the carrier is mostly inside the room: expectation,
prestige, national pride, and trained seeing. In cold fusion the carrier
leaves the room before the route is stable. The failure mode becomes faster,
more public, and harder to correct because the claim arrives in the public
mind with revolutionary consequences already attached.

The two cases together show why failure stories cannot serve as weapons
against curiosity. A culture that never risks weak signals will miss real
discoveries. A culture that never announces important possibilities will
hide work that matters. The question is how much social force the record has
earned. N-rays had not earned a field. Cold fusion had not earned a
revolution.

That proportionality is one of overhead's key disciplines: the carrier fits
the route. A lab notebook can carry a hint. A seminar can carry a promising
anomaly. A technical paper can carry a detailed claim. A press conference
carries a record strong enough to survive the public world it is about to
enter. When these scales are confused, overhead becomes an amplifier of
wish.

The N-ray scale mismatch was smaller than cold fusion's, but it was already
visible. Reports created enough literature and prestige that the claimed
phenomenon could seem more established than its route deserved. This is one
of the ordinary ways a record gets overcarried: not by one spectacular leap,
but by many small acts of deference. Each new report assumes a little too
much from the previous one, and by the time the structure is challenged,
the challenge feels like an attack on a community rather than a test of a
route.

The correction was therefore theatrical enough to break the theater. Quiet
skepticism might have been absorbed as lack of sensitivity. A public
interruption made the dependency visible. When overhead has become part of
the apparatus, correction can act on the overhead as well as on the
instrument. The experiment is no longer only optical; it is
social-optical.

The N-ray pattern --- a route too weak for the social load placed on it ---
appears at small scale before catastrophe. Mature propaganda resists
correction because it has already recruited identities, institutions, and
punishments. Early overhead can still be corrected by embarrassment,
interruption, and better controls. N-rays are a mercy in the history: a
public failure that ended before the carrier became an empire.

The mercy carries a duty. A good laboratory culture lets the apparatus
embarrass the observer, the observer embarrass the leader, and the critic
embarrass the group. If embarrassment is treated as disloyalty, the route
will be protected from the very event that could save it.

The episode therefore acts as an early gate. A community that makes room
for its own correction while the stakes are still small can absorb a larger
refusal later. Early severity is not hostility to discovery; it is mercy
toward the future field. That mercy is one of overhead's hidden virtues:
correction before identity.

\section{Cold Fusion: Announcement Before Record}
\episodecoord{1989 onward | Overhead :: REPEATABLE refusal -> PROPAGANDA}

Cold fusion belongs here because the press event became part of the
apparatus. The claim failed in later laboratories, but it also entered
public belief through a channel that outran the replication discipline.

In March 1989 the electrochemists Martin Fleischmann and Stanley Pons
announced at the University of Utah anomalous heat in electrochemical cells
with palladium electrodes immersed in heavy water. They interpreted the
heat as evidence of nuclear fusion occurring at room temperature. If true,
the result would have been extraordinary: cheap energy, a reconfigured
nuclear physics, and an engineering revolution arriving through chemistry
rather than reactors. That possibility gave the claim enormous public
force before the experimental route was ready to carry it.

The press conference changed the claim's order of operations. Instead of a
difficult result moving through detailed disclosure, critical replication,
and gradual public confidence, the announcement turned the claim into news
while many laboratories were still trying to learn what the Utah team
had actually carried out. The public heard possibility, controversy, and hope before
the record acquired a stable replication discipline. Overhead became one
of the forces acting on the experiment.

The underlying question was not foolish. Unexplained heat in an
electrochemical system is an experimental claim. It deserves calorimetry,
controls, materials analysis, nuclear byproduct checks, and replication.
The failure lies in the mismatch between evidential closure and public
force. A claim may be worth testing and still be unprepared to reorganize
belief. Cold fusion shows what happens when those two statuses collapse
together.

The apparatus was not a stage prop. Electrochemical cells, palladium,
heavy water, loading conditions, calorimetry, and heat balance all
mattered. Each of those components becomes a possible refuge when
replications fail. Perhaps the palladium was different. Perhaps the
loading was insufficient. Perhaps the calorimetry missed a boundary
condition. Perhaps the effect is real but rare. Such possibilities are
not automatically dishonest. They become overhead when they protect the
claim from ever facing a public halt condition.

Replication is the central refusal. Many laboratories attempted to
reproduce the effect under urgent public pressure. Some reported hints;
many failed; the field fragmented into dispute. The main record could
not stabilize around a repeatable nuclear effect strong enough to support
the announced interpretation. The cold-fusion claim could not merely
encounter skepticism; it encountered the ordinary scientific requirement
that other rooms be able to make the claim answer.

Extraordinary claims require no magical standards. They require ordinary
standards applied with enough severity for the claim's consequences. If a
system produces fusion, the heat story meets nuclear accounting:
nuclear products such as neutrons, tritium, or helium-4 of fusion origin;
energy balance; independent cell construction; and reproducible
protocols. When those links fail to hold, public excitement is not
allowed to fill the gaps.

A claim can develop an afterlife after the central record has failed.
Communities continue, specialized meetings form, new terms appear, and
believers distinguish themselves from hostile mainstream science. Some of
this persistence can be honest curiosity; anomalous measurements can be
worth revisiting. Persistence is not evidence by itself. A community's
survival may show that the hope survived rather than that the route
held.

The case is more subtle than ``press conference bad.'' Public
communication is not forbidden. Institutions have duties to report
important work, and secrecy can also corrupt a record. The problem is
premature public force. When a claim enters the world as revolution before
its methods can be tested, the social machinery begins creating winners,
losers, reputations, funding decisions, and ideological attachments
around a record that has not yet earned them.

Premature public force damages both believers and skeptics. Believers
inherit a claim already framed as world-changing, which makes correction
feel like suppression. Skeptics inherit a claim already framed as
overreach, which can make open anomalies feel like contamination. The
press event polarizes the experimental question. Instead of a question
about how a particular cell behaves under specified conditions, the
public hears ``are you for or against the revolution?''

This is how overhead degrades inquiry. It adds social identities faster
than the apparatus can add constraints. The cell becomes a flag.
Replication becomes loyalty or betrayal. A failed run becomes evidence of
incompetence or conspiracy rather than a data point in a method under
test. Once that happens, the experimental route is no longer the only
route by which the claim travels.

Hope is part of experimental life. People build hard apparatus because
they want the world to answer in new ways. Hope stays downstream of the
record. In cold fusion, hope became a carrier before the record had a
carrier. That inversion is the overhead failure.

The same structure now repeats in every premature-announcement culture. A
preprint, press release, corporate demo, viral graph, institutional
statement, health claim, machine-learning demo, materials announcement,
or crisis-science result can move faster than replication. The
technologies change; the structure is the same. Publicity creates load.
Once the claim has an audience, correction becomes harder. A result that
would have been a technical dispute inside a field becomes a public
drama with reputational stakes.

A second anchor lesson: the channel of distribution can become part of
the experiment's fate. A result announced too soon is not automatically
false, but it is placed under pressures that make falsity harder to admit
and truth harder to isolate. The public carrier scales to the evidential
carrier; otherwise the imbalance is overhead masquerading as discovery.

The calorimetry point deserves to be slow. Heat is a difficult record. It
can hide in gradients, losses, phase changes, recombination, calibration
drift, and boundary assumptions. A small unexplained excess in a
complicated cell is not the same kind of mark as a clear nuclear product.
If the interpretation is fusion, the heat meets nuclear accounting;
otherwise the word ``fusion'' carries more weight than the cell has
earned.

That mismatch is the exact place where announcement becomes dangerous. A
technical audience might hear ``unexplained heat in a difficult system''
and ask for cells, controls, blanks, loading ratios, neutron or tritium
checks, and calorimeter details. A public audience hears ``fusion in a
jar'' and rightly imagines civilization-scale consequences. The overhead
expands the claim's meaning before the apparatus has fixed the claim's
content.

Cold fusion damaged trust in a way that simple falsehood would not have.
After a revolutionary promise is announced and then refused by the
mainstream record, believers can interpret the refusal as institutional
suppression, and skeptics can interpret any continued work as
pathological. Both positions make careful anomaly work harder. The
premature carrier creates a social field in which ordinary experimental
patience looks politically loaded.

The lesson for institutions is uncomfortable. Universities, laboratories,
and funders want credit for important discoveries. They also carry a
duty to keep claims proportional to the record. That duty is not
public-relations timidity; it is evidential hygiene. The more
world-changing the claim, the more carefully the institution protects the
route from being overwhelmed by the announcement.

The structure repeats in fields unrelated to electrochemistry: health
claims, machine-learning demos, materials announcements, and crisis
science. A result with great promise enters a distribution channel
before uncertainty has been made durable. The channel generates belief,
investment, and identity. Later correction then fights not only an
experimental claim but the social world built during the interval.

The replication attempts are the moral center of the story. A replication
failure is not merely a negative vote; it is another laboratory trying to
stage the same claim under enough shared conditions that the result can
leave its origin. When many such attempts fail or produce unstable hints,
the public claim shrinks. It may shrink to ``there is an unresolved
calorimetric problem.'' It may shrink to ``some materials behave
strangely.'' It may shrink to ``the original interpretation was not
supported.'' It cannot remain world-revolutionary by borrowing energy
from the announcement.

Secrecy and priority pressure can also damage route sharing. If the exact
materials, preparation, and operating conditions are not available in
time for serious replication, the public claim has outpaced its method.
Patent concerns, institutional competition, and fear of being scooped are
understandable human pressures, and they are also overhead. When they
keep the route private while the claim is public, they invert the order
of experimental trust.

Cold fusion therefore sits between fraud and honest error in a more
interesting region. An overhead-failure verdict here is possible without
deciding that every actor was deceitful. Sincere researchers can
overinterpret. Institutions can overannounce. Audiences can overbelieve.
Critics can overpolarize. The failure is the whole arrangement by which a
fragile claim acquired a public body too large for its experimental
skeleton.

The public also learned the wrong kind of disappointment. A scientific
refusal teaches proportion: this record cannot support that claim. A
public debacle teaches cynicism: maybe all science is hype, or maybe all
rejection is suppression. Both are losses. Premature overhead risks more
than a false belief; it can damage the public's ability to understand
how ordinary refusal works.

Spectacular promises will continue to arrive. Energy, health,
intelligence, longevity, climate, and defense all generate claims whose
possible rewards dwarf the available record. The cold-fusion discipline
scales the announcement to the route. Hope motivates the work. Hope
cannot publish the conclusion.

\begin{phenom}{The Cold-Fusion Announcement Effect~\cite{fleischmann1989}}
\label{ph:cold-fusion-announcement}

\PhStatement
When announcement outruns replication, publicity becomes overhead attached to
an unclosed record.

\PhOrigin
Fleischmann and Pons announced anomalous heat in electrochemical experiments
interpreted as cold fusion. Subsequent attempts at replication produced dispute,
failure, and a long institutional afterlife.

\PhObservation
The public record included the laboratory claim, the press conference, the
institutional stakes, and the uneven replication attempts.

\PhConstraint
An extraordinary experimental claim cannot borrow certainty from the social
machinery that distributes it.

\PhConsequence
The refusal is not "new claims are bad." The refusal is that public force is
not a substitute for a repeatable record.
\end{phenom}

Cold fusion also reveals the afterlife of an unclosed record. Once the
mainstream verdict hardened, the continuing community defined itself
against that verdict. This is a common overhead pattern. A rejected claim
can survive as a technical research program, as a grievance identity, or
as both. The more the community's self-understanding depends on exclusion,
the harder it becomes to separate promising anomalies from the social
drive to prove that the exclusion was unjust.

A mirror risk runs on the mainstream side. After a public debacle, a field
may learn to dismiss adjacent anomalies too quickly because the name has
become embarrassing. That, too, is overhead. The record can refuse the
unsupported fusion claim without making electrochemical anomaly work
untouchable. Good science can say: that conclusion failed, that route
could not hold, and some narrower questions may still deserve method.

This balance is difficult because public culture prefers verdicts. Fraud
or breakthrough. Suppression or embarrassment. Genius or crank. Experiment
usually has more intermediate states: artifact, uncontrolled heat,
materials effect, rare condition, bad interpretation, suggestive but not
licensed, not yet repeatable, not worth the public claim. Overhead becomes
healthier when it has language for these intermediate states.

The cold-fusion story therefore teaches institutional proportion. A
university can be proud of risky work without spending the authority of
the institution ahead of the route. A journal can publish difficult claims
when the methods are inspectable without letting novelty substitute for
closure. A critic can demand replication without using embarrassment as a
substitute for reading the apparatus. Every role has an overhead virtue
and an overhead vice.

The claim's energy promise mattered because energy is never only a
technical subject. Energy claims carry geopolitical hope, economic
fantasy, environmental relief, military implication, and civilizational
anxiety. A new energy source is almost impossible to discuss as a modest
laboratory result. The laboratory route therefore deserves extra
protection. When the possible consequence is enormous, the social field
inflates the claim unless the institution holds it close to method.

Peer review and publication are not single magic gates. A paper can be
published and still require replication. A journal can make a claim
inspectable without making it settled. A conference can be legitimate
without being decisive. Overhead fails when these different public
statuses collapse into one binary: accepted or rejected, real or fake,
official or suppressed. Experimental records carry more levels than that.

Those levels are not bureaucratic niceties; they prevent words from
carrying too much. ``Reported'' is not ``replicated.'' ``Replicated'' is
not ``explained.'' ``Explained'' is not ``engineered.'' ``Engineered''
is not ``safe.'' ``Safe'' is not ``economically or politically adopted.''
Cold fusion compressed many such levels because the imagined payoff was
enormous. The discipline of overhead is to uncompress them again.

A short counterfactual: imagine the same anomalous heat claim introduced
with full methodological disclosure, no world-revolutionary press frame,
clear uncertainty language, and an organized replication program before
public conclusion. The result might still have failed as
fusion. The failure would have taught more and damaged less. The
difference is not the cell; it is the carrier.

After a field rejects an interpretation, the rejected community may
continue to produce technical claims under changed names. Some will be
weak. Some may be worth narrow investigation. The mainstream record needs
language for ``not fusion as announced'' without using the old failure to
pre-decide every later measurement. Otherwise the refusal itself becomes
overhead, a reputation field that can block proportionate curiosity.

The balance offers no rescue for the original announcement. It makes the
refusal more adult. A mature record can hold several sentences at once:
the announced interpretation failed; replication failed to support the
public claim; some materials and calorimetric questions may remain
technically interesting; none of that licenses the original revolution.
Overhead improves when it can speak in such differentiated sentences.

Renaming a contested program can clarify a narrower question, or it can
launder an old claim into a new public form. The test is whether the new
name carries stricter methods and humbler claims, or only escapes the old
embarrassment. After a debacle, names attempt to reset overhead.

Timing also matters. A claim can be premature in March and still become
testable later; or it can be suggestive in a paper and false
as a press conference. The same underlying measurement carries different
public meanings depending on when and how it is carried. Overhead is not
only the existence of publicity. It is the timing relation between
publicity and record.

The announced frame persists even in correction. Later reports are read
through the old promise; skeptics hear ``cold fusion'' before they hear
the method; believers hear ``suppression'' before they hear the failed
control. This is the half-life of overhead. A premature carrier decays
slowly, and its decay products remain active in the culture of the
field.

The cure is route-first language. Say what was measured before saying
what civilization might become. Say what failed before saying who was
vindicated. Say what remains open before building a community around the
opening.

Route-first language also protects the public. Citizens require no
calorimetric detail, but they benefit from knowing whether a claim is an
unreplicated anomaly, a replicated effect, a mechanism-supported theory,
or an engineered technology. Cold fusion gave too many listeners the
later statuses before the earlier ones had been earned.

When the institution calibrates these status levels, the public can
calibrate hope. When the institution collapses them, the public cannot.

\section{Lavoisier and Lysenko: Institution Is Not Evidence}
\episodecoord{1789 / 1948 | Overhead :: PROPAGANDA -> ACOLYTE}

An institution either carries evidence or replaces it. Lavoisier and
Lysenko answer that distinction by example. Lavoisier shows an institution
carrying a corrected experimental route into chemistry. Lysenko shows an
institution punishing the route that could have corrected it. The point is
not that politics lives outside science --- it never has --- but that an
institution can either preserve a record's refusals or recruit people to
repeat a claim after the record has failed.

Antoine-Laurent Lavoisier's chemistry was institutional in the strong
sense. Around 1789 he and his collaborators reorganized French chemistry
through laboratory apparatus, balances, nomenclature, pedagogy, and
publications. They were arguing against an older explanatory habit.
Phlogiston theory had treated combustion and calcination as the release
of an inflammable principle (``phlogiston'') from the burning substance;
metals on calcination were supposed to lose phlogiston into the air. The
trouble was that calcined metals weighed more than the original metals,
not less, and the gases involved had become amenable to collection and
measurement.

The new chemistry's public route was a disciplined accounting practice.
Substances entered and left reaction vessels under careful weighing;
gases were collected and identified; combustion was reinterpreted as
combination with oxygen rather than release of a hidden principle.
Conservation of mass disciplined what chemistry was allowed to claim had
happened. The reformed nomenclature --- oxygen, hydrogen, nitrogen, and
the systematic naming of acids, oxides, and salts --- made the new
accounting teachable. The balance was not only a tool; it was a public
grammar, and the new vocabulary pointed back to the balance.

That accounting refused phlogiston. Phlogiston had explanatory reach,
historical prestige, and teaching habits behind it. Lavoisier's program
made that convenience too expensive by forcing reactions into mass
balance, gas collection, and revised nomenclature. The institution
around the new chemistry carried the record's refusal into teaching,
terminology, and laboratory practice. The new chemical language became
overhead in the good sense: a support structure for a route that had
earned its authority.

The political biography is not clean. Lavoisier was also a tax farmer
for the old French monarchy and was executed by guillotine during the
Terror in 1794. The state killed a scientist whose
experimental work was reorganizing chemistry. That fact carries no claim
that science floats above politics, and no claim that politics can refute
a balance. The guillotine could not make phlogiston true. Political
violence can destroy a person without changing what the apparatus had
made public.

The distinction prevents a childish reading of institution. An
institution can be unjust in one register and still carry a correct
record in another. A person can be socially compromised and
scientifically decisive. The experimental gauge canonizes no persons; it
asks whether the record's route remains answerable. The balance and
accounting discipline survived Lavoisier's death because others could
use, criticize, and teach the route beyond his life.

Trofim Denisovich Lysenko worked under opposite institutional pressure.
From the late 1920s through the 1960s, his theories of plant breeding
--- including vernalization and the inheritance of acquired characters,
together with a hostility to chromosome-based heredity --- gained Soviet
state support, in part because they promised rapid agricultural
transformation under the urgent demands of Stalinist agricultural
policy. By 1948, with Stalin's backing, the Lenin All-Union Academy of
Agricultural Sciences declared Mendelian genetics ideologically
incompatible with Soviet biology. Geneticists were dismissed, imprisoned,
or executed; research programs were closed; teaching was rewritten. The
scientific problem was not simply that Lysenko had wrong ideas;
scientists have wrong ideas all the time. The institutional problem was
that the carrier of the claim gained power to punish the experimental
community that could refuse it.

Nikolai Ivanovich Vavilov names the lost route most concretely. Vavilov
ran the All-Union Institute of Plant Industry in Leningrad and led an
empirical program of crop diversity, centers of origin, plant collection,
and genetic reasoning. The institute's seed bank carried that program
forward as material memory: collections preserve variation so later
tests, breeding, comparisons, and recovery remain possible. The seeds
are not symbols. They are an apparatus for asking biological questions
across decades. Vavilov was arrested in 1940 and died in prison in 1943.
Lysenkoist institutions went beyond disagreement with a theory; they
damaged the conditions under which a record of biological variation could
continue to discipline agronomy.

The \lean{PROPAGANDA -> ACOLYTE} transition is severe here. Propaganda
is more than the announcement of a false sentence. It changes the
social cost of disagreeing with the sentence. An acolyte is then
produced: the person who repeats the institution's claim because the
institution controls reward, safety, status, publication, or survival.
The record no longer disciplines the institution; the institution
disciplines the record.

The hinge between the two cases is the word ``corrigible.'' A healthy
scientific institution can be wrong, and it remains correctable by
measurement, replication, and argument. It can have prestige, but
prestige loses when the route fails. It can have textbooks, but textbooks
change when the record changes. It can have leaders, but leaders remain
answerable to the apparatus. When those conditions disappear,
institutional success becomes a different predicate from experimental
success.

Lysenkoism also shows why propaganda prefers total explanations. A claim
that answers every difficulty by invoking ideological enemies,
insufficient commitment, or hostile conditions becomes hard to test,
because every failed test is reinterpreted as disloyalty. The route
disappears into the carrier. The claim survives by converting refusal
into moral defect.

The cases are not symmetrical in moral weight. Lavoisier's institutional
chemistry and Lysenko's state-backed biology are not equivalent uses of
authority. One let the balance and reaction ledger discipline chemistry.
The other let political authority discipline biology against genetics.
The point of comparing them is contrast. Institution is a carrier. The
carrier's virtue depends on whether it keeps carrying the route or
begins carrying itself.

The case also warns against a shallow anti-institutional lesson. Without
institutions, Lavoisier's chemistry would not have reorganized teaching
and practice; without seed banks and journals, agricultural genetics
would have had no durable memory. The answer to Lysenko is not no
institutions; it is institutions whose authority remains subordinate to
experimental refusal.

\lean{ACOLYTE} is funny only until it is not. In mild form, it is a
seminar room where students learn which sentences win approval. In
severe form, it is a state system where disagreement threatens
livelihood or life. The same structural test applies across the range:
can the acolyte still say what the record says when the record refuses
the institution?

A claim may be carried by institutions; it cannot be made true by them.
A carrier may amplify a record; it cannot punish the record's refusal.

Naming can be either liberation or capture. Lavoisier's nomenclature
made chemical accounting more public by forcing substances and reactions
into a shared grammar pointed at weighing. Lysenko's ideological
vocabulary made genetics suspect before experiment could speak. Names
are overhead because they carry claims between classrooms, papers,
ministries, and farms. A good name keeps the route reachable. A bad
institutional name makes the route socially dangerous to reach.

A student entering Lavoisier's reformed chemistry could learn a grammar
that pointed back toward balances and reactions. A student entering
Lysenkoist biology learned a grammar that pointed toward political
approval and away from genetic refusal. Education is never neutral
overhead. It trains future observers in what counts as an admissible
sentence. The health of an experimental culture depends on whether
teaching preserves the right to be corrected.

A seed collection is slow material memory. Destroying or marginalizing
it reduces the future's ability to ask biological questions; propaganda
can wound the future of measurement, not only the reputation of the
present. Institutions can starve a record, silence a community, redirect
training, and ruin lives before later correction arrives. Genetics
survived Lysenkoism; the intervening damage was real.

The positive lesson is equally serious. Institutions that preserve
refusal are among science's greatest inventions. A journal that
publishes a negative result, a collection that preserves awkward
variation, a standards lab that reports uncertainty, a university that
protects dissenting evidence, and a textbook that revises itself are
all forms of overhead serving the record. The enemy is not institution;
the enemy is institution made unanswerable.

The acolyte pattern in Lysenkoism was devastating because it recruited
agricultural urgency. Crop yields, famine, national survival, and
ideological modernization created pressure for usable answers, and a
claim that promises immediate transformation can outrun slow genetic
evidence. Propaganda often feeds on real need by offering a sentence
that cannot afford to be tested. Vavilov's empirical program represented
the slow route: collect, compare, preserve, test, breed, revise.
Lysenko's institutional success favored the fast sentence backed by
power. The fast sentence was wrong; power also made slowness look like
sabotage.

A crisis can make method look luxurious. Method is exactly what prevents
urgency from becoming fantasy. The more urgent the application, the
more dangerous it is to let institutions punish inconvenient evidence.
A field may need decisions before every question is settled; it cannot
abolish the questions in order to make decision feel settled.

The narrow modern warning is direct: whenever a carrier treats refusal
as moral or political deviation rather than as part of the record's
discipline, it has begun to move along the Lysenko vector. Modern
democratic assessment controversies are not Stalinist biology. The
severity differs; the structural vector is real.

Correct science is not socially weightless. The new chemistry needed
textbooks, demonstrations, laboratories, translations, allies, and
institutional prestige; a correct route still travels through people
and institutions. The difference is that the travel remained anchored
to operations that others could repeat. Naming oxygen carried weighing
forward rather than replacing it.

Lysenko's institutional vocabulary performed the opposite function: it
made the politically desired conclusion easier to repeat and the
experimental objection harder to say. Even honest observers may begin
adapting their reports to the language that keeps them safe. The
acolyte is not always a fanatic; sometimes the acolyte is a tired
person navigating a system that has made truth expensive.

Propaganda as overhead is cruel: it changes the cost structure around
observation. Ordinary evidence becomes dangerous; earned agreement
becomes prudence under pressure. The record no longer competes on
experimental grounds alone; it competes against a social system that
can punish its appearance. When a carrier takes ownership of truth,
human beings pay for the carrier's pride: people, collections,
students, crops, patients, and future questions all bear the cost.

Lavoisier's death sharpens this without producing a simple martyr
story. The state can kill a scientist and still fail to kill the route
if the route has entered public practice. The fact is terrifying and
hopeful at once: persons are vulnerable, and a well-staged record can
outlive the person who staged it. Institutions can destroy people; they
cannot vote a balance out of balance once enough others have learned to
use it.

Lysenkoism shows the darker mirror. Institutions cannot change the
underlying biology, but they can delay the public life of biological
truth and make the route harder to practice, teach, fund, or preserve.
Farmers, students, prisoners, and hungry people suffered while
propaganda governed the institutions whose duty was to carry the record.

Public usefulness is not evidence. Lavoisier's chemistry became useful
because its record held; Lysenko's claims promised usefulness under
ideological urgency. A promise of usefulness can be one of the most
powerful carriers of falsehood because it makes refusal look
anti-practical. A crop program that cannot tolerate genetics is not
practical; it is performing practicality.

The institution's memory is decisive. Lavoisier's chemistry left behind
practices that could be repeated by people who owed him no loyalty;
Lysenkoism tried to leave behind loyalty without leaving behind
repeatable practice. The contrast asks whether the carrier creates
independent competence or dependence on the carrier's approved
sentence.

Independent competence is the enemy of propaganda. A student who can
weigh, collect gas, grow crosses, inspect variation, or repeat an
analysis can be disciplined by the world. An acolyte trained only to
repeat the institution's answer is disciplined by reward. The
educational difference becomes a scientific difference.

The contrast therefore gives a standard for scientific education: teach
the route, not only the conclusion. A student who knows the conclusion
without the route can become an efficient repeater; a student who knows
the route can become a future corrector. Institutions that fear future
correctors have already confessed what kind of overhead they are. The
humane opposite of acolyte formation is teaching students to know where
trust comes from: trust earned by route can survive correction; trust
demanded by authority treats correction as threat.

\begin{phenom}{The Lavoisier--Lysenko Institution Effect~\cite{lavoisier1789,lysenko1948}}
\label{ph:lavoisier-lysenko-institution}

\PhStatement
Institutional success and experimental success are different predicates.

\PhOrigin
Lavoisier's chemical work used balance, conservation, and naming to refuse
phlogiston chemistry. Lysenkoism gained institutional authority in Soviet
biology while suppressing genetics and punishing dissenting experimental
communities.

\PhObservation
Both cases involve institutions, names, authority, and public teaching. Only
one lets the experiment continue to discipline the institution.

\PhConstraint
The social carrier of a claim remains answerable to the experimental record
it carries. When the carrier punishes the record's refusal, it has become
propaganda.

\PhConsequence
\lean{ACOLYTE} is the historical state in which a person repeats the system's
sentence because the system owns the reward, not because the record earned the
claim.
\end{phenom}

An LHC trigger is a finite rule for deciding which collision events get
kept. It may look like a relief after Lysenko --- morally it is --- but
structurally both cases involve institutions deciding what enters the
record. Lysenkoism decided by power against refusal. A trigger decides by
finite rule under public constraint. The difference is not that one has
selection and the other lacks it. The difference is whether selection can
be criticized by the record it shapes.

This distinction matters because modern science cannot avoid selection.
Fields select problems, instruments select events, journals select papers,
committees select priorities, and models select variables. The fantasy of
no selection belongs to no real science. The experimental discipline is
to make selection inspectable and corrigible. Lysenko hid selection behind
ideology. The LHC trigger names selection as method.

The contrast also prevents a lazy anti-politics. The LHC is political in
the ordinary sense that it requires states, budgets, treaties, agencies,
and collaborations. That political reality cannot turn its event selection
into propaganda. The question is whether political and institutional
structures remain answerable to detector behavior, calibration,
statistical review, and public method. Politics around a record is
inevitable. Politics replacing the record is the failure.

Lavoisier--Lysenko therefore frames honest finitude as something distinct
from coercion. Selection is not automatically suppression; authority is
not automatically propaganda. The gauge asks for the route of correction.
When a selected record can still say no to the institution, overhead has
not yet become acolyte machinery.

\section{The LHC Trigger: Honest Finite Selection}
\episodecoord{2000s onward | Overhead :: FINITE_ELEPHANT}

The Large Hadron Collider cannot keep everything. That fact is not scandal; it
is honest overhead. The trigger is a public confession that event streams are
larger than the record can carry.

Collider experiments produce events at rates far beyond what can be stored and
studied in full. The detector sees more than the archive can afford. This is
not a moral failure of particle physics. It is the physical and computational
condition of the experiment. The world gives too many collisions; the record
decides which ones become data. \lean{FINITE_ELEPHANT} is the name for
that condition: the elephant is real, large, and finite in what it can carry.

The trigger is therefore a staged loss. Hardware and software systems inspect
signals quickly, apply selection criteria, and reduce an overwhelming event
stream to a smaller population that can be stored for later reconstruction and
analysis. Some events are kept. Most are lost. The honest form of this loss
refuses to pretend the loss never happened; it publishes the selection
discipline by which loss became admissible.

This distinguishes the LHC trigger from dishonest forms of overhead. N-rays
lost control over expectation. Cold fusion lost balance between announcement
and replication. Lysenkoism punished refusal. The trigger
announces its limitation. It says: we cannot carry all events, so here is the
rule by which the finite record is made. That confession is not weakness. It
is the precondition for trust.

The trigger also makes theory and instrumentation meet under time pressure.
Selection rules are not arbitrary. They encode what the experiment is trying
to see, what backgrounds dominate, what detector signatures matter, and what
rates the system can tolerate. A trigger menu is therefore a public compromise
among physics goals, electronics, computing, storage, and uncertainty. It is
overhead because it is not the collision itself. It is evidence-supporting
overhead because it tells the record how the collision stream was narrowed.

The risk is obvious. A trigger can miss what nobody knew to ask for. A system
designed to keep certain signatures may discard unanticipated events. This is
not a trivial risk in a machine built partly to discover the unexpected.
Modern experiments answer by keeping multiple trigger paths, monitoring
efficiencies, preserving control samples, using calibration streams, and
designing searches with awareness of selection. The answer is not perfect
omniscience. The answer is visible finitude.

Visible finitude changes the status of a discovery. A Higgs candidate event,
a rare decay, or a missing-energy signature is not merely an event in nature.
It is an event that passed a chain of hardware, software, reconstruction, and
analysis choices. The final plot carries those choices inside it. The event
is not less real for having passed through them; it is scientifically usable
because the passage is documented.

This point is central to a realistic account of big science. Large
experiments are not just larger versions of small experiments. They move
some of the experimental intelligence into pipelines, collaboration
procedures, detector simulation, calibration databases, and trigger logic.
Imagining ``the data'' as a heap of raw collisions already misses the
experiment. The data are a finite public object made from an event stream
by staged loss.

Loss belongs in the account because it prevents selection from sounding
like triumph. Every finite record excludes possibilities. The trigger
becomes honest not by pretending no excluded event could have mattered,
but by giving later analysts a way to know which populations were
admitted, which were suppressed, and how selection efficiencies affect
the claim. The record's authority depends on that humility.

The trigger belongs in Overhead rather than Staging because it certainly
stages the event stream, but its main lesson is the cost of carrying a
record at scale. When collisions arrive faster than storage, overhead is
not optional. The experimental community builds a social and technical
apparatus for deciding what the public record can afford to remember.
That apparatus then answers to the scientific claims made from the memory
it creates.

The same structure appears outside particle physics. Telescopes filter
surveys. Genomics pipelines filter reads. Medical devices compress signals.
Climate archives homogenize station records. Social-science projects exclude
responses under prewritten rules. Every large record has a trigger-like
discipline somewhere. The LHC makes the discipline explicit because the event
rate is too obviously impossible to keep in full.

This explicitness is a virtue. A public rule for loss is better than a
hidden fantasy of completeness. The trigger teaches respect for finite
selection without worshiping it. It is neither scandal nor final
authority. It is a necessary compromise whose terms travel with the
result.

The trigger also reframes discovery as population management. A discovery
claim depends not only on the selected events but on the backgrounds and
efficiencies that selection created. The experiment knows how ordinary
processes populate the kept record before an excess can be interpreted as
new physics. Selection, simulation, and statistical analysis form one
overhead chain. The event that survives the trigger is only the beginning
of inference, not its end.

A trigger is costly, selective, collaborative, and theory-informed. It is
also honest when it lets the community see how the finite record came to
be. It proves that overhead can be good science when overhead confesses
its own shape.

The confession is quantitative. Trigger efficiencies, thresholds,
prescales, dead time, calibration samples, and simulated acceptances all
become part of the result's meaning. A selected event population is a
population under rules. If the rules are not known, background estimates
and signal claims float above the record.

The LHC also shows why independence can be designed inside a giant
system. Separate detectors, analysis teams, blinded regions, control
samples, and internal review processes create ways for the
collaboration to surprise itself. This is not the same as a lone
outsider repeating a tabletop experiment, but it answers the same
obligation under different scale conditions. Big overhead builds its
own internal refusals.

This internal refusal is easy to caricature as bureaucracy. Some of it
is bureaucracy. A detector collaboration without review and selection
discipline would not be more pure; it would be less answerable. The
relevant question is whether the bureaucracy protects the route or
protects the collaboration from the route. In the trigger case, the
public selection logic is a protection of the route.

The trigger is therefore a model for other finite records. The honest statement
is not "we kept everything." The honest statement is "we could not keep
everything; here is what we kept, why, with what efficiency, and what that
means for the claim." That sentence is one of the most mature forms of
experimental humility.

It also changes how to think about ``raw data.'' In a collider, rawness
is already a layered achievement. Electronics decide thresholds. Detector
channels are calibrated. Signals are synchronized. Events are built from
subdetector information. Reconstruction algorithms turn detector responses
into tracks, jets, leptons, missing transverse energy, and other objects. By
the time a physicist looks at an analysis dataset, the collision has become a
processed citizen of an experimental state.

This is not a scandal. It is how the collision becomes usable. The
scandal would be to hide the processing while invoking the romance of
raw nature. Big experiments make their audiences less sentimental about
data, not less trusting. Trust comes from documented transformation. A detector without
transformation is blind; a detector with hidden transformation is opaque; a
detector with public transformation is an instrument.

The finite trigger also forces a discipline of counterfactual thinking. What
would the experiment have seen if the signal were different? What backgrounds
pass the same selection? What rare but ordinary processes mimic the signature?
What classes of event were never admitted? These questions are not secondary
to discovery. They are the conditions under which discovery can mean anything.

There is a political humility here as well. Large collaborations decide
together what to keep. Those decisions are influenced by theory,
previous results, machine conditions, and computing limits. Honest
trigger culture cannot pretend these decisions came from nowhere. It documents them so later
claims can be interpreted under the finite memory that produced them. The
record is not omniscient. It is governed.

The LHC trigger thus gives a positive image of overhead under constraint.
Where N-rays hid expectation and cold fusion inflated announcement, the
trigger names finitude. It states what the experiment cannot accomplish.
That kind of limitation is one of the record's strongest forms of
honesty.

The trigger also creates obligations for later reinterpretation. Suppose a new
theory predicts events in a region the old trigger rarely kept. The archive
may not be able to answer. That is not betrayal if the limitation was public.
It is the finite record being finite. The honesty lies in letting future
researchers know where the absence of evidence is partly an absence created by
selection.

This matters because scientific archives often outlive their first theories.
Data collected for one purpose later become evidence for another. The
usefulness of that afterlife depends on metadata about how the record was
made. A kept event without trigger context is a fossil without provenance. It
may still be interesting, but its evidential force is damaged.

The LHC also teaches restraint in the opposite direction. Because the record
is finite and selected, one cannot demand impossible completeness before
allowing any inference. The question is not whether every event was kept. The
question is whether the kept population, selection model, and uncertainty are
adequate for the claim. This is the mature form of empirical reasoning under
load: neither omniscience nor cynicism.

The trigger hierarchy gives this maturity a concrete form. Fast hardware-level
decisions make crude but necessary reductions. Later software-level decisions
can be more refined. Offline reconstruction can apply still more detailed
calibration and analysis. Each layer has different information and different
time. The record is made by this stratified sequence. A single word,
"trigger," hides a layered economy of attention.

That economy also shows why simulations become part of overhead. To
understand what the trigger has admitted or rejected, experiments
simulate detector response, known processes, and possible signals.
Simulation is not the event, but it helps interpret the finite event
population. This creates another obligation: the simulation is validated
against control data. Otherwise the selection model can become a second
hidden apparatus.

The LHC's virtue is not that every outside observer can personally audit
every line of trigger code. The virtue is that the collaboration has
public norms for turning finite selection into documented uncertainty. The route is difficult,
but not mystical. This is how trust works at scale: not through personal
inspection by everyone, but through structured opportunities for competent
inspection by someone.

That form of trust is fragile and precious. It depends on a community willing
to publish enough detail for criticism, train enough people to understand the
detail, and preserve enough records for reinterpretation. A trigger hidden as
proprietary magic would not have the same epistemic status. The selection
would still happen, but the public record would be weaker because the loss
would be less answerable.

The LHC trigger therefore shows why complexity cannot excuse opacity.
Complexity may make full explanation difficult, but it increases the
duty to explain the parts that bear on the claim. The larger the
overhead, the more important it is to name where the public can test
the carrier.

This is also why open data, when possible, is not a slogan but a gradient.
Some data can be released broadly. Some require specialized formats, embargoes,
or safety constraints. Some require simulation and calibration files before
they mean anything. The honest question is not whether the word "open" appears,
but whether the record's usable route becomes more inspectable. Openness
without route is theater; route without any access is private authority.

The trigger's strongest lesson is therefore not technical but public: finite
records can be trusted when finitude is part of the record. The worst
selection is not selection itself. It is selection pretending to be
completeness.

The same lesson disciplines statistical summaries. A histogram, plot, or
significance value is a finite presentation of a selected population.
It carries enough account of the selection to make clear what the
figure can and cannot say. The trigger is a dramatic instance of a
universal rule: no selected record can pretend to be the unselected
world.

This rule is a bridge into every later compiled inference. A compilation is
made of selected records. If the selection is visible, compilation can be
honest. If selection disappears, compilation becomes authority without memory.

\begin{phenom}{The LHC Trigger Effect~\cite{lhctrigger2017}}
\label{ph:lhc-trigger}

\PhStatement
Finite selection is admissible when the rule for loss is part of the record.

\PhOrigin
Modern collider experiments use staged trigger systems to reduce enormous event
rates to data streams that can be stored and analyzed.

\PhObservation
Most events are discarded. The kept record is shaped by hardware triggers,
software filters, detector constraints, and physics priorities.

\PhConstraint
A finite record created by selection publishes the selection discipline that
made it finite.

\PhConsequence
The trigger is overhead, but not \lean{BULLSHIT}. It is the cost of making an
overwhelming event stream measurable.
\end{phenom}

The LHC trigger makes overhead's positive vocabulary possible. Until this
point, overhead has mostly appeared as distortion: expectation, announcement,
coercion. The trigger shows overhead as disciplined finitude. It lets the argument
say that losing information can be honest when loss is governed, measured, and
made part of the claim. This is a crucial correction to any naive ideal of
complete evidence.

No experimental record is complete in the metaphysical sense. A notebook
selects. A photograph frames. A trial endpoint chooses. A telescope points.
A survey samples. A climate model discretizes. A detector triggers. The
question is whether the selection is appropriate to the claim and visible
enough to be challenged. The LHC makes the selection problem impossible to
ignore because the event stream is obviously too large.

That honesty prepares the climate section. Climate assessment is also a
finite-selection problem, though in a different register. It selects datasets,
models, scenarios, confidence language, and summary statements from a vast
literature. The selection can be honest, but only if the boundary between
measured record, modeled projection, and policy implication remains visible.
The trigger's humility is exactly what a live consensus needs.

\section{Climate Consensus: Chain, Bound, and Policy}
\episodecoord{1980s onward | Overhead :: PROPAGANDA -> ACOLYTE}

Climate consensus is dangerous because it is live.
The gauge fails if it becomes a denial tract; it also fails if it
becomes a sermon. The positive task is to keep the measured chain
strong while keeping policy premises visible.

The measured chain is not empty. Greenhouse gases absorb and emit
infrared radiation. Atmospheric carbon dioxide has been measured
directly in the modern instrumental record and inferred through
ice-core air bubbles for earlier periods. Surface temperature records,
ocean heat content, cryosphere changes, sea-level records, and regional
observations form a broad physical account of a warming climate. The
IPCC's physical-science assessment states the attribution claim in
strong language: human influence has warmed the atmosphere, ocean, and
land. That factual chain remains intact under discomfort with policy
politics.

The chain is also heterogeneous. A thermometer reading is not the same
kind of record as a paleoclimate proxy. A measured CO\textsubscript{2}
concentration is not the same kind of record as a scenario-dependent
projection of late-century impacts. An attribution study is not the
same kind of object as a cost-benefit claim about one law. A climate
model is not the same kind of object as an ethical argument about
burden-sharing. A consensus that preserves those differences remains
scientific. A consensus that speaks as though all links carry the same
evidential texture begins to use overhead as authority.

Cheap symmetry --- ``both sides are political, so the record
disappears'' --- is refused. Measured warming remains measured.
Radiative physics remains physics. Attribution chains remain
substantial even where uncertainty is real. The opposite cheapness is
also refused: a measured chain cannot license every preferred action
as though policy choice were the last line of an experiment. The
record can authorize facts, risks, and bounds. It cannot by itself
choose discount rates, land-use priorities, energy tradeoffs,
geopolitical burdens, or tolerable levels of intervention.

\lean{PROPAGANDA} becomes structurally precise here. Propaganda is not
a strong claim --- science makes strong claims --- but false certainty
achieved by dropping the facts that would impose uncertainty. In one
direction, propaganda drops the uncertainty, values, and tradeoffs that
separate ``warming is real'' from ``therefore this exact policy is
scientifically mandated.'' In the other direction, it drops the
measured and attributed facts so that uncertainty in projections can be
used as cover for denying the record that already exists. Both moves
are refused.

The acolyte pattern appears when boundary questions become loyalty
tests. A person who distinguishes measurement from policy may be
treated as if they deny warming. A person who accepts warming but
questions a favored lever may be treated as morally suspect. A person
who emphasizes uncertainty may be using it honestly or smuggling
denial under the banner of skepticism. A person who emphasizes
consensus may be carrying a real record or smuggling policy under the
banner of science. The overhead problem is that public identity can
outrun these distinctions.

Climate communication inherits the ordinary burdens of scale. No one
reads every instrumental record, proxy paper, attribution analysis,
model intercomparison, and impact study. Assessment reports compile;
summary documents compress; journalists translate; governments
negotiate language; advocacy groups select examples; opponents select
errors; institutions attach logos and authority. All of this is
overhead. Some is necessary; without assessment and communication, the
record would not travel. Compression is also a dangerous act: it can
preserve a boundary or erase one.

The phrase ``the science says'' is both useful and dangerous. It is
useful when it points to a well-supported measured or attributed claim.
It is dangerous when it makes values disappear. Science can say that
emission pathways imply different warming ranges under specified
assumptions. It can say that certain interventions would affect
radiative forcing, air quality, land systems, economic sectors, or risk
distributions under modeled conditions. It cannot make a society's
final tradeoff morally automatic. Those decisions require politics,
ethics, engineering, law, and risk tolerance. Hiding that under ``science
says'' is overhead pretending to be measurement.

The research-priority question is a useful test. Emissions reduction is
a physically grounded lever because greenhouse-gas forcing is part of
the causal chain. It is not the only physical lever available for
discussion. Carbon removal, adaptation, land-surface change, albedo
modification, marine cloud brightening, stratospheric aerosol
concepts, and other interventions raise scientific and governance
questions. The National Academies' reflecting-sunlight report is
careful: such research is bounded, governed, and not a substitute for
mitigation. That caution is the point. A real scientific question may
be risky, politically fraught, or unready for deployment and still
deserve bounded research rather than ritual taboo.

Overhead becomes visible when a policy frame narrows curiosity before
the record has had its say. The narrowing can move in any direction. A
mitigation-only culture can underinvestigate feedback-side or
albedo-side possibilities because they disturb the moral architecture
of the preferred policy story. A geoengineering-enthusiast culture can
overstate technical readiness and hide termination shock, regional side
effects, governance risk, and moral hazard. A denial culture can invoke
those same risks selectively to block every action while refusing the
measured warming record. The gauge's job is to make the boundaries
public, not to baptize one tribe.

The word ``consensus'' needs the same care. In a healthy assessment,
consensus is not the replacement for evidence. It is the public summary
of a large evidential field after disagreement has been worked through.
That kind of consensus is valuable because no individual can personally
audit every link. Consensus becomes propaganda when it ends questions
that the record itself leaves open. A consensus statement can summarize
the physical basis; it cannot make every downstream choice look like
heresy-proof entailment.

Climate belongs in Overhead rather than only in compiled inference
because the carrier and the inference have different jobs. The carrier
is the object. Assessment reports, public campaigns,
international conferences, activist rituals, industry counter-campaigns,
model ensembles, and policy packages all carry climate claims. Some
carriers preserve the measured chain. Some distort it. Some amplify
warranted risk. Some convert risk into identity. Some convert
uncertainty into denial. The record can be separated from the carrier
without pretending it travels alone.

The live case also teaches humility about language. ``Settled'' can
mean the core physical and attribution claims are strong enough to
guide public reasoning. It can also be abused to mean no further
boundary questions are legitimate. ``Uncertain'' can mean quantified
ranges, scenario dependence, and model spread. It can also be abused to
mean no factual inference is warranted. Both words can be honest; both
can become overhead. The gauge asks what each word performs in the
sentence.

Climate consensus therefore gives overhead its most contemporary
acolyte test. Can a person say all of the following in the same room:
measured warming is real; human influence is central; projections carry
scenario and model uncertainty; policy decisions include values;
mitigation, adaptation, removal, and reflecting-sunlight research each
have different evidence and governance profiles; uncertainty licenses
neither denial nor automatic command? If the room cannot tolerate that
full sentence, overhead has begun to govern the record.

The refusal is not against climate science. It is against the
collapse of climate science into a public authority blob that cannot
state where measurement ends and decision begins. A strong record can
survive such distinctions. A propaganda system cannot. Making the
boundaries visible therefore protects the measured chain from both
denial and overclaim.

Assessment and advocacy ride different sentences. If the assessment
voice speaks as advocacy without naming the move, it spends scientific
authority on a political choice. If the advocacy voice pretends the
evidence is weaker than it is, it spends uncertainty on delay. The
honest form names the register.

The scenario problem is another boundary. Scenarios are not
predictions in the same sense as tomorrow's measured temperature
forecast. They are conditional stories: if emissions, land use,
technology, policy, population, economic growth, and other drivers
follow such paths, then models project such ranges. The conditional
stories are useful. They become misleading when the conditions are
hidden and the range is presented as fate. They become equally
misleading when their conditionality is used to pretend no risk has
been shown.

Proxy reconstructions require the same honesty. Ice cores, tree rings,
sediments, corals, boreholes, and other archives can carry climate
information from before instrumental records, and each proxy has
translation costs. Those costs leave the proxies usable as scientific
objects. The reconstruction earns trust by carrying calibration,
uncertainty, and cross-checks. A public story that treats every
ancient climate number as if it were read from a modern thermometer
has lost the route.

Policy packaging also tends to hide distribution. ``Action is needed''
can conceal who pays, who benefits, who bears land-use changes, who
loses reliability, who gains new industry, who faces imposed
adaptation, and who gets protected from risk. These are not details the
climate record can dissolve. They are the public decisions that come
after the record. Calling them science because science supplied part of
the risk account is a category error.

Uncertainty is not a moral escape hatch. The thickness of the policy
problem cannot license denial of the measured chain. The atmosphere
remains measurable when a carbon tax is controversial. Ocean heat
remains ocean heat when integrated assessment models contain value
judgments. Denial is the mirror propaganda of overclaim: it achieves
false comfort by dropping facts that impose risk.

The double refusal preserves the record's force while forcing public
language to stop borrowing more force than the record gives. The
climate case is not a side fight. It is the contemporary place where
Overhead proves it can make distinctions under pressure.

The albedo question illustrates the needed distinction. Earth's
reflectivity is a physical part of the climate system shaped by snow,
ice, clouds, aerosols, vegetation, land use, and ocean surfaces.
Proposed interventions that alter reflection are not thereby safe,
governable, ready, or morally equivalent to emission reduction. They
are physical levers whose research status belongs in public argument
rather than in taboo or enthusiasm.

Climate institutions also handle the asymmetry between global and
local claims. Global mean temperature can be a meaningful metric while
local effects differ widely. A policy justified by global benefit can
impose local costs; a local objection can be real without refuting the
global record. Overhead becomes propaganda when either side erases
scale: when global risk erases local burden, or when local
inconvenience erases global evidence.

The same scale problem appears in time. Past warming is measured.
Future warming is projected under scenarios. Damages and adaptations
unfold through social futures. The farther the claim moves from
measured past to policy future, the more conditional it becomes. Good
public reasoning handles conditionality; propaganda hates it.

Short slogans are often overhead devices. ``Believe science'' may
defend a real record against bad-faith denial, and it can also blur
the distinction between physical assessment and political program.
``Climate has always changed'' may point to real paleoclimate
variation, and it can also erase the contemporary forcing record. The
question is not which slogan wins. It is which route each slogan is
hiding.

The climate case tests whether a community can preserve multiple truths
under institutional pressure. Warming is real. Human influence is
central. Impacts and risks matter. Uncertainty ranges matter. Policy
values matter. Engineering feasibility matters. Distributional justice
matters. Research portfolios matter. None of these facts cancels the
others. The overhead failure is the public machinery that needs one of
them to disappear so its preferred sentence can sound complete.

The IPCC process is a particularly complicated carrier because it is
both scientific and intergovernmental in its public form. Scientific
assessment, expert judgment, calibrated uncertainty language,
governmental review, and summary negotiation all meet in the final
documents. This complexity leaves the documents valuable rather than
worthless. Their authority has a route, and the route benefits from
being explained rather than mystified. A sentence in a summary has
passed through a carrier whose job is to preserve the assessed record
rather than to make the sentence politically frictionless.

Dissent needs the same care. Some dissent is bad-faith denial funded or
amplified to protect interests. Some dissent is real scientific
argument about sensitivity, feedbacks, datasets, regional projections,
damage functions, engineering feasibility, or policy design. A healthy
overhead system can tell the difference without pretending the
distinction is always easy. An unhealthy system treats all dissent as
denial or all consensus as corruption. Both moves save thought by
sacrificing the route.

The climate claim is therefore narrow but demanding: keep the chain
intact. Let measured warming remain measured. Let attribution remain
attribution. Let scenarios remain conditional. Let policy remain
policy. Let uncertainty remain quantified. Let action arguments state
their values. A strong record survives a hard politics this way, and
a citizen can then see which parts of the public claim are record,
which are projection, and which are choice.

Public confidence language --- likely, very likely, high confidence,
low confidence --- carries uncertainty without collapsing it into
vagueness. These terms can sound bureaucratic; they are also a
discipline. When public retellings drop them, they often drop the
feature that makes an assessment scientific.

Risk communication adds another layer. A low-probability high-damage
outcome may matter differently from a high-confidence moderate
outcome. Science can map such risks; it cannot decide a society's risk
tolerance by itself. Overhead becomes honest when it names whether a
sentence describes probability, severity, exposure, vulnerability,
feasibility, or value judgment.

The distinction between mitigation and adaptation is another boundary.
Reducing future forcing, preparing for impacts, removing carbon,
protecting ecosystems, hardening infrastructure, and researching
controversial interventions are not the same kind of response. Each
has evidence, costs, uncertainties, and moral questions. A public
carrier that treats one response as the only scientific one has
smuggled a value structure into the word ``scientific.'' A public
carrier that treats the existence of tradeoffs as a reason for inaction
has made the opposite smuggle.

Policy resolution comes after grammar: facts with their uncertainties,
models with their scenarios, risks with their distributions,
interventions with their side effects, and decisions with their values.
Climate consensus becomes scientifically stronger when that grammar
stays visible.

The same grammar protects scientists. When public language makes every
policy claim sound like a direct deliverance of science, scientists
inherit responsibility for decisions they could not make as
scientists. When public language lets opponents deny the measured chain
because they dislike policy, scientists are forced to relitigate facts
that already form the starting point. Boundary discipline is therefore
not only fairness to the public; it is protection for the scientific
record itself.

The final diagnostic is whether the consensus can name what would
change its mind. A scientific consensus remains scientific when new
measurements, improved models, or better attribution studies could
revise its claims. A policy identity may resist such revision because
revision threatens the identity. The measured warming record is strong
enough to survive revision in details. A propaganda carrier is not.

This criterion makes revision difficult rather than easy. Climate
records are complex and slow. Some uncertainties narrow over decades,
not months. Some disputes are resolved by new observing systems;
others by model comparison; others by waiting for outcomes no ethical
person wants to test by inaction. Difficulty leaves intact the
obligation to say what kinds of evidence would alter the claim. A
consensus that cannot describe its own possible correction has become
a creed.

Correction is not reversal. A revised estimate of sensitivity,
sea-level range, aerosol forcing, or regional risk is not necessarily
a reversal of the warming record. It may be the record becoming more
exact. Propaganda wants every correction to be either hidden or
weaponized. Science treats correction as maintenance.

That maintenance is public because climate decisions are public. A
private technical distinction is not enough when the public sentence
carries political force. Assessment language travels with enough
structure that citizens can see which parts of the claim are record,
which are projection, and which are choice.

The climate case also makes visible the difference between expertise
and mandate. Experts can say what follows from radiative physics,
observations, models, and risk analysis. They can also, as citizens,
argue for policies. The two roles may be held by the same person, and
the public sentence marks the transition. If the transition is hidden,
expertise is spent as mandate. If the transition is marked, expertise
can inform decision without pretending to replace it.

Marking that boundary makes bad-faith criticism less effective. Denial
thrives when it can point to overclaim and pretend the overclaim
refutes the record. Keeping the boundary clean deprives denial of that
gift. The measured chain carries only what it has earned. Policy
advocates need not be shamed for advocacy; the requirement is that
action arguments state their additional premises --- risk tolerance,
equity, cost, feasibility, moral priority --- so that compelling
premises are recognized as premises rather than as thermometer
readings.

Climate science is strong enough to live with visible boundaries. It
matters because the record is real, the risks are real, and the
decisions are public. The public record can force responsibility
without pretending that responsibility is itself a measurement. That is the adult form of consensus: strong facts,
explicit uncertainty, and openly owned decisions.

A brittle consensus defends every public use of its authority because
any boundary looks like weakness. A strong consensus affords boundaries
because its core record is unburdened by overclaim. It can say what it
knows, how it knows it, what follows conditionally, and where society
decides.

The experimental record is strongest when it is not asked to pretend
that politics vanished. It can enter politics as a disciplined
constraint without entering politics disguised as total command. The
same discipline protects legitimate urgency. Urgency that names its
values can be argued with. Urgency that hides inside measurement
becomes brittle. Climate needs the former, because the record is
already heavy enough without being asked to impersonate every later
choice.

\begin{phenom}{The Climate-Consensus Boundary Effect~\cite{arrhenius1896,ipcc2021,nas2021reflecting}}
\label{ph:climate-consensus-boundary}

\PhStatement
A compiled scientific consensus preserves the boundary between measured
facts, quantified uncertainty, model projection, and licensed action.

\PhOrigin
Modern climate assessment reports compile instrumental records, paleoclimate
proxies, radiative physics, attribution studies, and model ensembles into a
public scientific account.

\PhObservation
Some links are close to direct measurement, such as atmospheric carbon dioxide
records and temperature series. Other links are chained through proxies,
models, scenarios, economics, engineering constraints, and ethical choices.

\PhConstraint
The consensus cannot smuggle policy conclusion into measurement, and critics
cannot use uncertainty in projections to override the facts that the record
already licenses.

\PhConsequence
The gauge refuses both failure modes: certainty produced by dropping
uncertainty, and denial produced by pretending uncertainty erases fact.
\end{phenom}

Leaving climate behind erases nothing. The climate record stands; its later
inference questions stand; the urgency of the political consequences stands.
Overhead simply marks the boundary at which a live public consensus carries
immediate political weight, and shows why such a consensus deserves unusually
strong boundary discipline. The more urgent the consequences, the stronger
the temptation to collapse the chain. Boundary discipline answers that
temptation.

Climate also forces Overhead to admit that communication is itself an
experiment in public trust. Too little force, and a real risk is ignored. Too
much compression, and the public is asked to obey a sentence whose internal
boundaries have disappeared. Too much caveat, and bad-faith denial exploits
the caveat. Too little caveat, and honest uncertainty is treated as sabotage.
No slogan solves this. Only disciplined speech solves it.

Disciplined speech can say "measured," "inferred," "projected,"
"scenario-dependent," "value-laden," and "decided" without shame. Those words
sort the record rather than weaken it. The public deserves to know which
parts of the chain are thermometers, which are proxies, which are models,
which are ethical choices, and which are institutional preferences. The
sorting is the opposite of denial. It is the condition for legitimate action
under a real record.

Overhead reaches farther than any one political battlefield. Observatories,
citations, replication, registries, samples, fossils, and control charts
show the same structure in quieter form. The live case tests the gauge under
pressure; the ordinary cases keep the gauge honest.

\section{Big Observatories: Load as Architecture}
\episodecoord{1900s-2000s | Overhead :: LOAD}

Big observatories make \lean{LOAD} visible without making it shameful. A
large telescope, interferometer, satellite mission, neutrino detector, or sky
survey is not merely an instrument. It is architecture, schedule, maintenance,
software, calibration, access policy, weather, labor, archive, and institutional
competition. The record exists because the load is carried.

This is one reason romantic histories of experiment become misleading. A
single image from a space telescope or a single event from a detector can look
like a direct gift from nature. Behind it are proposals, review panels, launch
systems, detectors, cooling systems, operations teams, calibration campaigns,
pipeline software, data releases, and failure-management plans. The mark is
public because a whole overhead structure made the mark possible.

The danger is not that such structures exist. The danger is that their scale
can become persuasive in itself. A billion-dollar instrument can tempt the
public into thinking the answer is important because the machine is
important. It can tempt the collaboration into protecting the machine's
justification. It can tempt critics into treating all big science as theater.
All three temptations are wrong. Load is neither evidence nor fraud. It is the
burden that remains answerable to evidence.

Big observatories also change authorship. A result may list hundreds or
thousands of collaborators. The old picture of a discoverer at a bench no
longer fits. Authorship becomes a record of infrastructure as much as
interpretation. That can be honest: the detector, calibration, software, and
operations all contributed to the result. It can also make responsibility
diffuse. If everyone is an author, who owns the exact route from source to
claim? The answer is procedural: internal review, public documentation,
data products, calibration notes, and reproducible analysis.

The archive is part of the instrument. Modern observatories often produce
records that outlive the first question they were built to answer. A sky
survey can support later discoveries because the data were preserved with
metadata and calibration. That archival load is real science. It lets the
record travel beyond the original team. It also creates a new risk: later
users can detach a data product from the pipeline conditions that made it.
Good archives carry provenance, not just files.

Big observatories therefore teach a balanced respect. They are not temples
whose cost proves their truth. They are not bureaucratic monsters whose cost
discredits their truth. They are load-bearing structures. The scientific
question is whether the load remains visible, whether the selection and
calibration routes are documented, whether outsiders can challenge the claim,
and whether the instrument's institutional survival is kept separate from any
particular result.

The maintenance story is often where the real science hides. Mirrors degrade,
detectors age, electronics drift, cryogens boil, software dependencies change,
ground stations fail, and calibration sources need their own custody. A
beautiful observation can depend on a thousand unglamorous acts of keeping the
instrument within its intended state. That maintenance is not separate from
the data. It is what allows data taken in different months, years, or
instrument modes to be compared.

Large instruments also create access politics. Who gets observing time? Which
proposals count as important? Which data are proprietary, and for how long?
Which teams get early access to pipelines and calibration knowledge? These
questions cannot make the science unreal, but they shape which questions the
instrument answers first. \lean{LOAD} includes the governance of attention.
The sky may be common, but telescope time is not.

The history of big science also shows why public trust requires more than
spectacular images. Images travel easily; calibration papers travel slowly.
The public may see a nebula, a black-hole shadow, or a gravitational-wave
chirp, while the evidence rests on processing choices and instrument models.
Good overhead builds paths from spectacle back to method. Bad overhead lets
spectacle become proof.

A big observatory's archive can be more democratic than its construction. Few
people can build the instrument. More people can later interrogate the data if
the archive is documented and accessible. The archive cannot erase inequality
of expertise, but it lets the route outlive the first gatekeepers. A public
record grows stronger when later questions can arrive after the original
prestige has cooled.

The load also changes failure. When a tabletop experiment fails, one may lose
time, pride, and materials. When a satellite, telescope, or detector fails,
the loss is distributed across careers, agencies, national promises, and
future observing opportunities. That distributed cost can make institutions
less willing to hear bad news. A failing component, flawed calibration, or
software error becomes not only technical trouble but a threat to a public
investment story. Good big science builds channels for bad news before the
bad news arrives.

Engineering reports, anomaly reviews, calibration memos, and operations logs
belong in the history of experiment. They are not secondary paperwork. They are
the record of load meeting reality. A mission that admits
instrument drift, failed modes, pointing errors, or pipeline revisions may
look less heroic, but it is more scientifically usable. The route has joints,
and the joints are documented.

Big observatories also teach that public access is not the same as public
understanding. Releasing data without releasing calibration context can
produce a new kind of false openness. The file is available, but the route is
not. Honest openness includes documentation, uncertainty, processing history,
and warnings about known limitations. Otherwise access becomes spectacle in a
different costume.

The best large instruments therefore have two public faces. One face shows
what the instrument found. The other shows how the instrument knows. The first
face attracts attention; the second preserves trust. Overhead becomes
dangerous when the first face is funded and polished while the second is left
to specialists no one is expected to consult.

This two-face structure is unavoidable. Public science needs images, stories,
and discoveries that can be communicated. It also needs calibration,
uncertainty, and operations detail that few nonspecialists will read. The
failure is not that one face is more accessible. The failure is when the
accessible face is allowed to stand as if the technical face were absent.

Big observatories therefore require a culture of translation that refuses to
collapse into simplification by erasure. The public can be told the result plainly
while still being told that a route exists: an archive, a pipeline, a
calibration paper, an uncertainty statement, a collaboration review. Trust is
not produced by making every citizen an expert. It is produced by making the
expert route publicly real.

This is the most charitable account of institutional authority. The public
often trusts expert communities because the route is too difficult to walk
alone. But honest trust points toward the route, not away from it. "Trust
us" is weaker than "here is how this can be checked by those equipped to check
it, and here is what remains uncertain." Big observatories are at their best
when they teach that distinction.

This is also how large instruments earn democratic patience. A public may not
understand every calibration, but it can understand that calibration exists,
that error is reported, that data have provenance, and that the institution
refuses to let cost stand in for truth. That is not perfect transparency. It
is accountable opacity.

Accountable opacity may be the best phrase for much modern experiment. The
record is too technical for total public transparency, but it need not be
opaque in the authoritarian sense. It can expose enough structure for trust to
have handles: who checked, what was selected, where the data live, what
uncertainty remains, and how correction would appear.

Those handles are what separate expertise from priesthood. The public cannot
perform every calibration, but it can see that calibration is not a private
sacrament. The route has doors, even if only trained people can walk through
all of them.

\begin{phenom}{The Big-Observatory Load Effect~\cite{galison1997}}
\label{ph:big-observatory-load}

\PhStatement
Large instruments create necessary overhead whose authority remains
subordinate to the records they stage and carry.

\PhOrigin
Twentieth- and twenty-first-century big science reorganized experimental work
around large collaborations, major instruments, pipelines, archives, and
institutional review.

\PhObservation
The record includes apparatus, operations, calibration, collaboration
procedures, software, data releases, and access rules.

\PhConstraint
Cost, scale, and collaboration size may support a record only by preserving
the route to the claim. They may not substitute for the route.

\PhConsequence
\lean{LOAD} is honest when it carries evidence and dangerous when it begins to
ask evidence to justify the carrier.
\end{phenom}

Machines give way to texts and fields. Big instruments make load visible
because concrete objects are built and maintained. Literature load is less
visible. A citation network or replication culture has no single dome or
detector, but it carries claims just as surely. The overhead becomes habits
of reading, reviewing, repeating, and rewarding.

This hidden load can be more influential than architecture because it enters
almost every scientific sentence. A telescope may be used by a limited
community; citation habits are used by everyone who writes. If those habits
preserve origins, they make knowledge portable. If they preserve fluency
without origins, acolytes form without anyone issuing explicit orders. The
carrier becomes style.

\section{Citation Cascade: Acolytes Without Malice}
\episodecoord{modern | Overhead :: ACOLYTE}

A citation can be honest overhead. It lets a claim travel without forcing
every later scholar to rebuild the entire route. But a citation can also
become an acolyte mechanism. A paper cites another paper, which cites another paper,
which cites a weak or misread source. The claim acquires a trail of scholarly
surface area without acquiring new evidence.

The failure is usually not dramatic fraud. It is often ordinary compression.
Authors summarize. Reviewers accept familiar sentences. Introductions repeat
field lore. A claim becomes background. Eventually the sentence is not being
carried because someone has inspected the route; it is being carried because
the literature has learned to say it. This is \lean{ACOLYTE} in a mild but
pervasive form: repetition as membership in a discourse.

Citation cascades are especially dangerous because they imitate traceability.
A later analyst sees many references and assumes the claim has many supports. But a
cascade may have only one root, or the references may support adjacent claims,
or the original result may have been qualified in ways the cascade has
removed. The form of scholarship remains while the route to the fact becomes
blurred.

None of this means citation is corrupt. Citation is one of the great
inventions of public knowledge. It is the way claims travel with enough
address for later analysts to find their origin. The problem is address
without custody. Honest citation makes a claim more inspectable. When it
makes the claim merely more respectable, overhead has begun to replace
evidence.

The cure is not to demand that every sentence contain a full replication. The
cure is to keep important claims connected to their evidential roots, to
distinguish review from primary result, to mark uncertainty when a citation is
used for background, and to make correction acceptable when the cascade has
overstated its root. A living literature keeps the power to prune inherited
sentences.

Citation cascades also create a particular kind of courage problem. Correcting
an inherited sentence can feel petty because the sentence is already part of
the field's scenery. The person who asks for the root may be treated as
slowing things down, missing the point, or failing to know the literature. In
that moment, the literature has become an institution protecting its own
fluency. The acolyte is not bowing before a tyrant; the acolyte is preserving
smoothness.

The smoothness can matter in high-stakes fields. A background claim about a
disease mechanism, a risk factor, a material property, or a historical event
may guide future experiments and funding. If the cited root is weak, the
field may spend years optimizing around an inherited convenience. The cost is
not just a bad sentence. It is misdirected attention.

Good citation practice is therefore not scholastic fussiness. It is overhead
discipline. Honest citation tells a later analyst where a claim can be
wounded. If the cited source is a review, say so. If the primary result is
small or contested, carry that condition. If the claim is field lore, leave
it bare rather than decorating it with references that never actually bore
its weight. The route matters even in prose.

Citation cascades also interact with prestige. A claim from a famous lab,
journal, or theorist can enter the literature with more momentum than its
evidence alone would deserve. Later authors may cite the prestige because
others have already cited it. The cascade then turns original status into
apparent consensus. This is not a conspiracy; it is inertia wearing a
bibliography.

The problem is intensified by review articles. A good review can clarify a
field and direct later inquiry to roots. A bad review can launder
uncertainty into background. Once a review sentence becomes the source everyone cites, the
primary record may stop being read. The field then inherits the review's
compression as if it were the experiment's own voice.

Citation also has a negative version: noncitation. Failed replications,
qualified results, corrections, and null findings may be omitted from the
lineage because they make the story harder to tell. The visible cascade then
looks cleaner than the field. Overhead is not only what gets carried. It is
also what gets left behind.

The practical virtue is citation with friction. Important claims deserve a
pause long enough to let the writer ask whether the cited source really
supports the sentence. That pause is small, but it is one of the places where a
literature remains answerable. A field that loses the pause becomes fluent at
its own expense.

The pause also has a social cost. It can slow writing, irritate reviewers, and
make a familiar introduction less smooth. That is why citation discipline is
overhead rather than mere virtue. It requires time and sometimes awkwardness.
A healthy literature values that awkwardness enough to keep inherited
sentences from becoming untouchable.

A citation cascade is easiest to correct early. Once a sentence has entered
textbooks, grant language, and review articles, correction feels disruptive.
The source route therefore stays visible from the beginning. The
longer a weak claim travels, the more overhead gathers around the act of
admitting it was weak.

The practice of retraction and correction belongs here too. A retraction is
not only the removal of a paper. It is a public signal that the carrier has
allowed the record to interrupt it. Retractions can be messy, politicized, and
late, but a literature that cannot retract has no immune system. Citation
discipline includes knowing when a route has been closed.

The healthiest citation culture is therefore not one without error. It is one
where error can be followed, named, and stopped. A literature that can say
"this line no longer supports that claim" has preserved its route.

That sentence is hard to say because citation is social memory. Correcting it
changes what a field remembers about itself. But memory that cannot be edited
is not scholarship. It is liturgy.

\begin{phenom}{The Citation-Cascade Acolyte Effect~\cite{greenberg2009}}
\label{ph:citation-cascade}

\PhStatement
A claim can acquire scholarly force by repetition even when the evidential
root has not grown stronger.

\PhOrigin
Studies of citation distortion and citation bias show how biomedical and
scientific claims can travel through selective or inaccurate citation networks.

\PhObservation
The record includes cited claims, secondary summaries, review articles,
selective quotation, and field lore.

\PhConstraint
Honest citation preserves inspectability. The number of carriers is not the
same as the strength of the root evidence.

\PhConsequence
\lean{ACOLYTE} need not be malicious. It can be the ordinary scholar repeating
what the literature rewards as already known.
\end{phenom}

Citation cascade and replication crisis belong together because one concerns how
claims travel in prose and the other how they travel in experiment. A citation
can make a claim look settled by repetition. A replication can unsettle the
claim by asking the world to answer again. The two practices need each other.
Citation without replication can become lore. Replication without citation can
forget what it is testing.

The good literature is therefore not the literature with the most references.
It is the literature whose references keep routes open for correction. That
means the literature tolerates embarrassment. A field can say, "we repeated
this sentence too easily," just as an experimental group can say, "we
selected this event population too confidently." The correction is not
humiliation. It is maintenance.

\section{Replication Crisis: Population-Scale Refusal}
\episodecoord{modern | Overhead :: REPEATABLE refusal}

The replication crisis is overhead turning back on itself. Journals,
incentives, small samples, flexible analysis, publication bias, and novelty
pressure helped carry many results farther than their records could support.
Then large-scale replication projects made the refusal public. The important
lesson is not that psychology or biomedicine is fake. The lesson is that
fields can develop overhead that rewards publishable claim-shapes before it
rewards durable routes.

Replication is a public act of humility. It says that a result earns its
weight only by surviving another route, another team, another sample, or
another analysis. This humility can feel threatening in
fields where careers and theories have grown around original findings. But the
threat is a virtue. A record that cannot tolerate replication pressure is not
being protected; it is being insulated.

The crisis also clarifies why "statistically significant" was never enough by
itself. A small p-value can be produced inside a fragile practice. If many
hypotheses are tried, measures chosen flexibly, exclusions adjusted after
looking, and null results left unpublished, the visible literature will
overstate the strength of effects. The number has the costume of measurement
while the overhead has shaped what numbers become visible.

Large replication efforts cannot produce final truth either. They have their
own design choices, sample differences, fidelity questions, and interpretation
debates. But they change the field's posture. They make repeatability a public
object rather than an assumed virtue. They force methods, materials, effect
sizes, power, and analysis plans into view. That is overhead becoming
corrective.

The replication crisis therefore belongs with both \lean{BULLSHIT} and
\lean{LOAD}, but its main name is refusal. A population of attempted repeats
can tell a literature that its carrying machinery has overcarried. The answer
is not despair. The answer is preregistration, registered reports, open data
where ethical, better power, clearer analysis plans, and less worship of
novelty. Those reforms are overhead too. They are overhead designed to make
overhead answerable.

The emotional difficulty is real. A failed replication can feel like an attack
on a scientist's competence or integrity, even when the replication is simply
asking whether an effect is stable. This emotional charge is part of overhead.
Careers, grants, labs, graduate students, and theories gather around results.
When a result fails, more than a number is threatened. The public record is
built so that this threat cannot decide the evidential question.

The crisis also exposes the weakness of single-study thinking. A lone study
can be interesting, suggestive, and methodologically careful while still being
too small to carry a general claim. A literature becomes trustworthy by
learning how effects behave across designs, samples, contexts, and analytic
choices. Replication is not clerical duplication. It is how the field learns
which parts of a result are stable and which belonged to the original setting.

This makes effect size as important as yes-or-no significance. A result may
replicate in direction but shrink in magnitude. Another may fail under a more
diverse sample. Another may depend on exact materials. These outcomes are not
all the same. A mature replication culture gives the field a vocabulary for
partial survival rather than forcing every result into triumph or fraud.

The replication crisis is therefore not an anti-science episode. It is one of
science's self-repair episodes. The refusal came from within the culture of
experiment and statistics. It asked whether the carriers of credibility had
become too generous. That question is painful because it is loyal to the
record.

The reforms are not free. Preregistration takes time. Larger samples cost
money. Registered reports require journals to change incentives. Open data
requires curation and ethical judgment. Replication work is less glamorous
than novelty. All of these are overhead burdens. The point is that they are
burdens designed to make the record less vulnerable to appetite.

The crisis also reveals a conflict between discovery culture and measurement
culture. Discovery culture rewards surprise, novelty, and clean narrative.
Measurement culture rewards stability, error accounting, and repeatability.
Science needs both, but the publication system can overpay discovery culture
until measurement culture becomes an afterthought. Replication projects are a
way of making measurement culture visible again.

The language of failure also carries weight. A failed replication is not
always a verdict that the original result was fraudulent or worthless. It may mean the
effect is smaller, context-dependent, method-sensitive, or absent. These
distinctions matter because a field that treats every nonreplication as
scandal will hide failures; a field that treats every nonreplication as a
minor nuisance will learn nothing. Refusal needs vocabulary.

Population-scale refusal therefore matures a field by giving it a distribution
of outcomes. Some effects survive. Some shrink. Some depend on conditions.
Some disappear. The field learns not one verdict but a map of reliability.
That map is overhead at its best: organized memory of where claims held and
where they failed.

The map also changes incentives for young researchers. If a field rewards only
novel positive findings, it teaches students to hunt for publishable surprise.
If it rewards careful replication, null results, and method improvement, it
teaches students that truth maintenance is real work. The replication crisis
therefore becomes a training crisis. What kind of scientist emerges from the
overhead?

This question links directly to \lean{ACOLYTE}. A bad incentive structure can
create acolytes of novelty without anyone intending fraud. A better structure
can create practitioners who expect their claims to be retested. The
difference is not personality. It is the ecology of reward around the record.

The crisis also changes the meaning of expertise. An expert is not someone
whose published effects are never questioned. An expert is someone trained to
let effects face the procedures that might reduce or remove them. That is a
harder identity, but a more scientific one. It makes correction part of
competence rather than humiliation.

That harder identity is exactly what overhead can cultivate. The carrier
earns trust by rewarding people who make their claims easier to test, not
only easier to admire. Replication reform is cultural engineering in service of humility.

Humility here is not meekness. It is technical courage: the willingness to
build a record that can reduce one's favorite effect. A field that teaches
that courage has begun to repair its overhead.

\begin{phenom}{The Replication-Crisis Refusal Effect~\cite{openscience2015}}
\label{ph:replication-crisis}

\PhStatement
A field's publication machinery can overcarry claims until organized
replication makes the refusal visible at population scale.

\PhOrigin
Large collaborative replication projects in psychology and related fields
tested whether published findings could be reproduced under new attempts.

\PhObservation
The record includes original effects, replication protocols, effect sizes,
power, analysis plans, successes, failures, and interpretation disputes.

\PhConstraint
A literature may not treat publication as a substitute for repeatability.
Public claim strength answers to replication pressure.

\PhConsequence
The crisis is a correction of overhead by overhead: institutions of repeat
testing are built to discipline institutions of publication.
\end{phenom}

The replication crisis also prepares the trial registry because it exposes a
missing-denominator problem. A visible literature is not the same as the
population of attempts. If only successes are visible, the field cannot know
how surprising success is. If only clean narratives are visible, the field
cannot know how many analytic paths were tried. Replication repairs some of
that after publication. Registration tries to repair it before publication.

This is an important shift in time. Replication says: let the result face
another run. Registration says: let the run be public before the result
exists. Both are overhead against selective memory. One looks backward at a
literature; the other protects the future literature from being born selected.

\section{Clinical-Trial Registries: Load Against Selective Memory}
\episodecoord{modern | Overhead :: LOAD -> LOGICAL seed}

Clinical-trial registries are ordinary and severe. They exist because an
unregistered trial can disappear too easily. If a study is planned, run, and
then left unpublished because its result is inconvenient, the public
literature becomes a selected memory. The treatment may look better than the
record warrants because the missing trials have no public address.

Registration adds load before the result is known. It records the trial,
intervention, endpoints, eligibility criteria, and planned analysis in a place
others can later inspect. This is bureaucratic overhead in the best sense. It
slows the claim down so that desire cannot rewrite the trial after seeing the
outcome. It also makes absence visible. A registered but unpublished trial can
be asked about.

The registry is not evidence of efficacy. A bad trial is not improved by
registration. Registration cannot prevent every analysis trick. Its job is
narrower and essential: it keeps the route from vanishing. The clinical record becomes more
public because the study's existence and planned shape are no longer owned
only by the sponsor or investigator.

Treatment decisions depend not only on positive findings but on the
population of trials that could have spoken. Selective memory is an overhead failure; registration is overhead
designed against it. The public record needs a ledger of attempts, not merely
a gallery of successes.

Registries also discipline endpoints. If a trial changes its primary outcome
after seeing the data, the change may be scientifically explainable, but it
stays visible. Otherwise the published paper can present a post-hoc success
as if it were the planned test. The registry leaves later analysts a way to
ask which claim was promised before the result and which claim was discovered
afterward. Both may matter, but they carry different inferential status.

The registry is also a public memory for harms and nulls. Patients accepted
risk by entering the trial. Their participation cannot vanish merely because
the outcome failed to flatter the sponsor, investigator, or theory. This gives
registration an ethical force beyond method. It honors the subjects by making
their contribution part of the public record even when the result disappoints.

This is overhead at its most quietly humane. Forms, numbers, dates, endpoints,
and status fields can look bureaucratic, but they protect inference from a
very human temptation: tell the story that came out best. The registry says
the story began before the ending was known.

Registries also protect meta-analysis. A later review that combines trials
depends on the denominator of attempts. Without registries, the review may
see only the bright surface of published success. With registries, reviewers
can look for missing results, changed endpoints, early stops, or abandoned
studies. The public inference becomes less dependent on sponsor storytelling.

This matters because clinical decisions convert evidence into action quickly.
A weakly supported treatment can harm patients, waste resources, or crowd out
better care. A beneficial treatment can be delayed if evidence is hidden or
fragmented. Registries cannot solve every clinical problem, but they make the
evidential field harder to curate in secret.

The registry's deeper lesson is that overhead can be preventative. Many cases
of overhead failure surface only after inflation. Registration tries to
prevent one kind of inflation by making the attempt public before the outcome
is known. It is a fence built before the cliff, not a rescue after the fall.

Preventative overhead is often unpopular because it looks like delay before
there is a scandal to justify it. Forms, protocols, and public identifiers
feel heavy when everyone is confident. They feel necessary after the missing
trials are discovered. A mature scientific culture funds the prevention before
the embarrassment, because the embarrassment costs patients and trust.

The registry is therefore a forward-looking promise: this study had a public
shape before it had a flattering ending. That promise is simple enough to look
administrative and deep enough to change inference.

\begin{phenom}{The Trial-Registry Load Effect~\cite{dickersin1992,icmje2004}}
\label{ph:trial-registry-load}

\PhStatement
Registration is honest overhead when it prevents inconvenient experiments from
disappearing before inference begins.

\PhOrigin
Concerns about publication bias and selective reporting led to stronger norms
and requirements for registering clinical trials.

\PhObservation
The record includes trial identity, endpoints, eligibility, planned analysis,
sponsor, status, and reported or missing results.

\PhConstraint
Clinical inference depends on knowing which trials existed, not only which
favorable trials reached publication.

\PhConsequence
\lean{LOAD} becomes logical preparation: a public ledger protects later
decision from selective memory.
\end{phenom}

Polywater, Piltdown, and Shewhart charts are smaller cases, not decorative
ones. They show three ordinary ways overhead gets corrected: sample custody,
renewed testing, and control of variation. They keep the analysis from
sounding as if overhead were only a matter of media, ideology, or giant
machines. The everyday record is always at risk of being overnamed,
overdesired, or overinterpreted.

Clinical trials make the ethical stakes explicit; the smaller cases make
the discipline portable. A water sample, a fossil fragment, and a factory
process all ask the same question in smaller form: what has the carrier
added, and can the route still refuse it?

\section{Polywater: Chemistry Refuses an Inflated Claim}
\episodecoord{1960s | Overhead :: BULLSHIT refusal}

Polywater was a claimed anomalous form of water with unusual properties. For a
time it attracted serious attention, because water is central and the reported
properties were strange enough to matter. Its afterlife depended on
contamination, enthusiasm, and the difficulty of keeping ordinary material
discipline visible once an extraordinary substance had been named.

The refusal came through chemical custody. If the anomalous properties could
be explained by impurities, dissolved substances, biological contamination, or
handling, then the new water was not a new form of water. It was a failure of
source discipline. The glamorous noun had outrun the sample.

Polywater matters because the corrective act is so ordinary.
Analyze the sample. Ask what is in it. Check the vessel, the capillary, the
handling, the environment. Chemistry asked no metaphysical argument against
polywater. It asked the claim to answer to material controls. The
inflation collapsed when custody became stricter.

The case also shows how naming accelerates overhead. Once "polywater" exists
as a noun, researchers and the wider public can imagine an object with a
stable identity. The name gets ahead of the custody. A sample still owed
suspicion becomes a member of a proposed kind. The corrective discipline
is to make the noun answer to the sample route.

The case is small, but the rule is large: a name cannot stabilize before
the source has. Once a name circulates, it collects
explanations, hopes, and opponents. Chemistry's answer was admirably plain.
Before asking what polywater means, ask what is in the water.

That plainness is the discipline. Overhead often inflates by abstraction. Custody deflates by returning to the vessel, the impurity, the
handling, and the sample. The route gets smaller and truer.

The same move applies whenever an impressive name outruns a controlled
source. Return to the source. Ask what entered the vessel, the archive, the
model, or the trial. The fastest way to deflate bad overhead is often to make
the material path embarrassingly concrete again.

\begin{phenom}{The Polywater Custody Effect~\cite{franks1970}}
\label{ph:polywater-custody}

\PhStatement
An extraordinary material claim fails when ordinary sample custody explains
the anomaly better than the proposed new substance.

\PhOrigin
Polywater claims surfaced in the 1960s and were later rejected when
contamination and experimental artifact became clear.

\PhObservation
The record includes anomalous water properties, sample preparation, vessels,
impurities, and chemical analysis.

\PhConstraint
A new material identity earns the authority of a named substance only after
surviving custody checks.

\PhConsequence
\lean{BULLSHIT} is refused by ordinary chemistry when the sample route becomes
more persuasive than the inflated name.
\end{phenom}

Polywater leads to Piltdown because both cases show names outrunning custody.
One name promised a new material. The other promised an
ancestral form. In both cases, desire condensed around an object before the
object's route had been made severe enough. The corrective work was not
glamorous. It was the slow reassertion of material facts over a name that had
become too attractive.

\section{Piltdown Man: Desire as Institutional Carrier}
\episodecoord{1912-1953 | Overhead :: PROPAGANDA}

Piltdown Man is not an experiment in the narrow sense, but it is an overhead
case in the history of evidence. Fossil fragments were taken as a major human
evolution find because they fit desires about ancestry, geography, and the
shape of human development. Museums, experts, and national prestige helped
carry the claim until later analysis exposed the composite fraud.

The important lesson is not simply that fraud happens. Fraud happens in every
human enterprise. The deeper lesson is that desire can make institutions slow
to apply the tests they already know how to value. A record that flatters a
community's self-image receives a smoother path. The carrier wants the claim
to be real, and the wanting affects the speed and severity of refusal.

Piltdown also shows why physical evidence needs continuing custody. A fossil
is not self-interpreting. Provenance, anatomy, chemistry, dating, comparison,
and access all matter. If the evidence is guarded by prestige instead of
reopened by method, the institution has become a shelter for falsehood. Later
testing accomplished more than correcting a date. It broke an institutional spell.

The long delay matters. A false record leaves damage while it lives. It
shapes textbooks, research expectations, and public imagination. Later
exposure is essential, but it cannot give back the decades spent reasoning
around a desired object. Overhead can waste historical time. That cost
belongs in the verdict.

Piltdown is therefore a warning about fit. Evidence that fits a preferred
story deserves more scrutiny, not less. The more welcome the object, the
more necessary the custody. Desire is not a reason to reject evidence, but it
is a reason to test the carrier.

The institutional shame is not only that the object was false. It is that the
object was allowed to be comfortable. Mature science grows nervous when a
specimen flatters the room too well.

That nervousness is not paranoia. It is respect for desire's ability to become
evidence-shaped. Piltdown teaches that a comfortable object needs an
uncomfortable test.

\begin{phenom}{The Piltdown Institutional-Desire Effect~\cite{weiner1955}}
\label{ph:piltdown-institutional-desire}

\PhStatement
Institutional desire can carry a false record when the claim flatters the
community authorized to test it.

\PhOrigin
Piltdown Man was accepted for decades as an important human ancestor before
chemical and anatomical tests exposed it as a fraudulent composite.

\PhObservation
The record includes fossil fragments, provenance claims, museum authority,
comparative anatomy, chemical testing, and delayed refusal.

\PhConstraint
Prestige and fit with expectation may not shield physical evidence from
renewed testing.

\PhConsequence
\lean{PROPAGANDA} can be quiet: not a shouted slogan, but an institution
letting a desired object keep speaking after tests were ready to interrupt it.
\end{phenom}

Piltdown leads to Shewhart because both involve institutional attention over
time. Piltdown shows attention captured by a desired object. Shewhart shows
attention disciplined by a chart. The contrast is useful:
overhead can teach people what to notice, but it can also teach them when not
to invent a story. Good control practice is institutionalized skepticism.

\section{Shewhart Charts: Ordinary Control Against Theater}
\episodecoord{1920s | Overhead :: LOAD discipline}

Shewhart control charts are humble overhead. They announce no new particle
and no new cosmos. They ask whether a process is behaving within expected
variation or whether an assignable cause has entered. This is the ordinary
discipline that many spectacular histories omit: before a claim can be
dramatic, a process learns what ordinary variation looks like.

The chart matters because it separates noise from action. If every fluctuation
produces intervention, the process becomes worse. If real shifts are dismissed
as mere noise, the process drifts. Statistical control is therefore a moral
discipline for measurement systems: ordinary scatter calls for patience;
structured departure calls for response. The residue discipline becomes
industrial overhead.

This is honest \lean{LOAD}. Someone collects the measurements, plots the
limits, maintains the process, and trains users not to overreact or
underreact.
The chart is not evidence of product quality by magic. It is a practice that
keeps variation visible enough for decisions to be governed. It resists
theater because it makes ordinary fluctuation boring and assignable causes
specific.

The chart is also an anti-acolyte device. It tells supervisors, workers, and
inspectors not to treat every story as a cause. The line either crosses a
rule, or the pattern earns investigation. This protects the process from charisma.
Someone may have a favorite explanation, but the chart asks whether the
process has actually departed. It is a small republic of variation.

That small republic matters because organizations love stories. A manager can
blame a worker, a worker can blame a machine, a team can blame a supplier, and
a consultant can sell a narrative. The control chart preserves judgment but
slows storytelling until variation has had its say. It is overhead
against narrative appetite.

In that sense Shewhart is overhead's quiet hero. The chart gives no grand
cosmology. It gives a way to keep ordinary processes from being narrated into
panic or complacency. Honest overhead often looks exactly that plain.

The plainness is the point. A process that knows its ordinary variation is
harder to manipulate with urgency, blame, or wish. Control is overhead that
has learned to distrust drama.

\begin{phenom}{The Shewhart Control-Load Effect~\cite{shewhart1931}}
\label{ph:shewhart-control-load}

\PhStatement
Control overhead is honest when it distinguishes ordinary variation from a
process change that deserves action.

\PhOrigin
Shewhart's statistical quality-control work introduced control charts for
monitoring industrial processes.

\PhObservation
The record includes repeated measurements, control limits, process history,
ordinary variation, and assignable causes.

\PhConstraint
A process record cannot convert every fluctuation into alarm or every
departure into noise.

\PhConsequence
\lean{LOAD} disciplines attention by making routine variation public before
theater can claim it as discovery or crisis.
\end{phenom}

Shewhart closes the ordinary sequence because control is the modest virtue
that Overhead needs most. The cases have shown spectacular failures of
carrying, but most overhead failures begin as small failures of control:
forgetting the baseline, changing the endpoint, trusting the familiar
sentence, protecting the desired object, or reacting to every fluctuation as
news. Control charts are not the whole answer. They are a visible symbol of
the answer: keep variation public enough that appetite cannot narrate it
freely.

\begin{movementclose}
Overhead closes by separating necessary cost from false authority. After that
separation, truth still travels. A claim that stays true only in the room
where it was made is not yet public science.

The cases have moved through several kinds of carrier. N-rays showed
expectation carrying a weak signal. Cold fusion showed announcement carrying
a claim before replication could. Lavoisier and Lysenko showed institutions
carrying records in opposite ways: one by preserving the balance's refusal,
the other by punishing biological refusal. The LHC trigger showed finite
selection carrying an overwhelming event stream honestly by naming the loss.
Climate consensus showed assessment, communication, and policy carrying a
real measured chain while constantly risking boundary collapse. Big
observatories carry load as architecture. Citations carry claims through
literature, with or without custody. Replication projects make overcarried
claims answerable. Trial registries keep inconvenient absences from
vanishing. Polywater, Piltdown, and Shewhart charts show that ordinary
sample custody, renewed testing, and statistical control can refuse theater
before theater becomes truth-shaped.

The first lesson is generous: truth travels with carriers.
Records decay without institutions. Standards drift without maintenance.
Large event streams vanish without triggers. Patients disappear from the
literature without registries. Telescopes and detectors run on cooling,
calibration, and crew, not on purity. The social and technical machinery
around experiment is part of how the universe becomes public.

The second lesson is severe: carriers are not self-validating. Prestige is
not signal. Announcement is not replication. Consensus is not a policy
proof. Citation count is not root evidence. Cost is not truth. A community
can carry a record, or it can begin to carry itself by using the record as
emblem. The same forms that make science possible can make imitation of
science possible.

The test of an honest carrier can be stated simply: the carrier remains
under correction by the record, or it has become authority theater. That
test applies to a dim screen in Nancy, an electrochemical cell in Utah, a
Soviet biology classroom, a collider trigger, a climate assessment, a
clinical registry, a museum fossil, and a factory chart. The scale changes;
the obligation holds.

Overhead also clarifies why public science is never merely a private
relation between mind and world. The record moves through institutions
without being swallowed by them. It accepts help without selling its
refusal. It compresses knowledge for use while keeping the route
recoverable. It trains acolytes in the good sense --- people disciplined by
the record --- while resisting acolytes in the bad sense, people rewarded
for repeating the carrier's sentence.

That tension is the hinge into Truth Transport. Carriers are not in
question; their existence is the precondition for travel. The question
becomes how a claim travels through them without losing witness, locality,
uncertainty, and refusal. A truth that cannot travel remains provincial. A
truth that travels by dropping its route becomes propaganda. Truth Transport
asks how experimental truth can move without becoming placeless.

Travel is also the moment when carriers become most tempting. The farther a
claim moves from its origin, the more people rely on summaries, authorities,
labels, and institutions. The route compresses or it cannot travel; but it
cannot compress so far that local conditions disappear. Truth Transport
therefore asks for transported locality: witness that can move, conditions
that can be reconstructed, and universal claims that remain answerable to
the local routes that earned them.

Overhead has also prepared a distinction between public truth and public
volume. A claim may be loud because it is well supported, or loud because the
carrier is powerful. A claim may be quiet because it is weak, or quiet because
the carrier has not yet learned how to move it. The danger is confusing
loudness with truth.

Every public claim now earns its hearing by showing both its carrier and its
route. If it can show both, the conversation continues. If it shows only
carrier, the failure has a name. If it shows only route but cannot travel,
Truth Transport has a task.

The carrier-route question also makes overhead retrospective. Entry made
records public. Residue made leftovers durable. Staging made values
portable. Overhead asked what happens when portability becomes valuable
enough to attract systems of protection, amplification, and reward. Every
earlier success becomes a risk here. A public record can become a slogan. A
durable residue can become a sect. A portable value can become a brand. The
cure is not retreat from publicity, durability, or portability. The cure is
keeping their routes visible under use.

The honest examples --- LHC trigger, trial registries, Shewhart charts,
archives, careful assessment language --- are not side notes. They prove
that overhead can learn humility. A carrier can name its losses, keep a
ledger before the outcome, distinguish ordinary variation from cause,
preserve provenance, and publish uncertainty. Such forms are not glamorous,
but they belong to the conditions for public truth.

The dishonest or failed examples --- N-rays, cold fusion, Lysenkoism,
citation cascades, polywater, Piltdown --- prove the mirror. They show what
happens when the carrier asks for more trust than the route can bear.
Sometimes the excess is gentle. Sometimes it is spectacular. Sometimes it
is deadly. The structure is the same: the record is no longer free to
discipline the sentence being carried.

Custody binds the procession. Under custody, a mark becomes entry; a
leftover becomes residue; a value becomes staged scale. Overhead is
custody under social pressure. Truth Transport is custody under distance.
Decision is custody under action. Experiment has never been left behind;
the path simply follows experiment into the forms that let it become
public.

The believer-skeptic shortcut closes here as well. Believers cannot cite the
carrier as proof. Skeptics cannot cite the carrier as disproof. Institutions
cannot ask trust to replace route. Critics cannot pretend routes travel
without institutions. The only durable path is slower: follow the claim,
inspect the carrier, recover the route, and ask what remains answerable.

That slower path is the realism of an experimental history. The universe
became public through carriers that were necessary and fallible. To
understand experimental history at all is to understand both halves.

Only after that audit can the large words be used again. \lean{SCIENTIFIC},
\lean{TRUTH}, \lean{WITNESSED}, \lean{REAL}, \lean{LOCAL}, and
\lean{UNIVERSAL} are dangerous words after Overhead. They have earned their
danger. Their use is allowed only when Truth Transport keeps them attached to the
conditions that let them travel.

The result is not suspicion of truth. It is suspicion on behalf of truth.
Overhead has shown how easily the social forms of knowledge can become
substitutes for knowledge. The opposite possibility --- that social and
technical forms can carry a local record outward without emptying it ---
depends on custody.

Overhead therefore ends without cynicism. The carriers failed in some cases
and served in others. The failures are warnings; the successes are
instructions. Public truth is not born by escaping overhead. It is born by
forcing overhead to remain answerable.

That answerability rests on a refusal of purity. There is no pure record
outside all institutions, and there is no institution pure enough to
replace the record. Public science lives between those impossibilities. It
builds carriers, audits them, repairs them, and sometimes tears them down.
The hinge from public cost to public truth is therefore not purity. It is
disciplined transport: a way for local records to cross rooms,
laboratories, nations, instruments, and generations while keeping enough of
their origin to remain answerable. Without this audit, travel would sound
too easy. With it, travel becomes an earned condition: not the
disappearance of cost, but the successful custody of a claim through cost.

The rule that survives is austere and generous at once: trust the carrier
in proportion to its willingness to be corrected by the route. That rule
refuses propaganda and sustains real institutions. It lets science have
laboratories, journals, assessments, archives, and public authority without
letting any of them become substitutes for experiment. It also refuses
despair, because every bad carrier has an answering good form: expectation
answered by blinding, announcement by replication, coercion by corrigible
institution, finite loss by declared selection, consensus theater by
boundary discipline, citation cascade by source custody, publication bias
by registration, narrative appetite by control. Overhead is dangerous
because it is powerful. It is useful for the same reason.

Overhead has therefore added one more obligation to the procession. A
public claim is not only marked, held, staged, and scaled. It survives its
carriers. It passes through the institutions that make it usable without
surrendering its power to correct those institutions. That obligation is
hard, but it is the only way experimental truth becomes more than a local
achievement.

Truth Transport begins at exactly that edge: the claim is carried, but the
route is not allowed to vanish.
\end{movementclose}
